\documentclass[10pt,numbers]{sigplanconf}
\pdfpagewidth=8.5in 
\pdfpageheight=11in 
\usepackage[final]{oopsla2016}

% The following oopsla2016 options are available:
%
% 1stsubmission   For the initial submission
% 2ndsubmission   For the 2nd submission
% final           For camera-ready

\usepackage{amsmath,amssymb,amsopn,amsthm}
\usepackage{mathtools}
%\usepackage{float}
\usepackage[utf8]{inputenc}
\usepackage[T1]{fontenc}
\usepackage{algorithmicx,algorithm}
\usepackage[noend]{algpseudocode}
\usepackage{multirow}
\usepackage{cleveref}
\usepackage{graphicx}
\usepackage[usenames,dvipsnames]{xcolor}
\usepackage{verbatim}
\usepackage{listings}
\usepackage[inline]{enumitem}
\usepackage{relsize}
\usepackage{xspace}
\usepackage{fix-cm}  % this is for math inside \scriptsize text
\usepackage{wrapfig}
\usepackage{mychangebar}
\usepackage{microtype}
%\usepackage[firstpage]{draftwatermark}
%\SetWatermarkText{{\hspace*{8in}\includegraphics[scale=0.1]{img/aec-badge}}}
%\SetWatermarkAngle{0}
\newif\ifdiff

%%%  diff hacks  %%%
\ifdiff
% configure changebar
\setlength\changebarsep{1mm}
% configure marginpar
\setlength\marginparsep{2mm}
\setlength\marginparwidth{1.5cm}
\newcommand\diffnote[1]{\hspace{0pt}\marginpar{\hbadness=10000\small #1}}
\newcommand\diffnoteup[1]{\marginpar{\hbadness=10000\small #1}}
\else
\nochangebars
\newcommand\diffnote[1]{}
\newcommand\diffnoteup[1]{}
\fi

%\makeatletter 
%\def\@cb@firstcolumntrue{}
%\def\@cb@firstcolumnfalse{}
%\makeatother

\usepackage{tikz}
\usetikzlibrary{fit,positioning,patterns,shapes,shapes.multipart,calc}
\usepackage{pgfplots}
\usepackage{pgfplotstable}
\pgfplotsset{compat=1.12} 

\include{macros}

\newif\iftr
 \clubpenalty = 10000
 \widowpenalty = 10000
 \displaywidowpenalty = 10000
\begin{document}

\copyrightdata{978-1-nnnn-nnnn-n/yy/mm} 
%\doi{nnnnnnn.nnnnnnn}

% Uncomment one of the following two, if you are not going for the 
% traditional copyright transfer agreement.

%\exclusivelicense                % ACM gets exclusive license to publish, 
                                  % you retain copyright

%\permissiontopublish             % ACM gets nonexclusive license to publish
                                  % (paid open-access papers, 
                                  % short abstracts)

\title{Deriving Divide-and-Conquer Dynamic Programming Algorithms using 
  Solver-Aided Transformations}

\authorinfo{Shachar Itzhaky\and 
            Rohit Singh\and
            \makebox[6cm][c]{Armando Solar-Lezama}\and 
            Kuat Yessenov\and 
            Yongquan Lu\and
            \makebox[6cm][c]{Charles Leiserson}}
           {MIT, USA}
           {\vspace{-2em}}
\authorinfo{Rezaul Chowdhury}
           {Stony Brook University, NY, USA}
           {}

\toappear{}

\maketitle

\begin{abstract}
We introduce a framework allowing domain experts
to manipulate computational terms
in the interest of deriving better, more efficient implementations.
It employs deductive reasoning to generate provably correct efficient implementations from a very high-level specification of an algorithm,
and inductive constraint-based synthesis to improve automation.
Semantic information is encoded into program terms through the use of refinement types.

In this paper, we develop the technique in the context of a system called Bellmania that uses solver-aided tactics to derive parallel divide-and-conquer implementations of dynamic programming algorithms that have better locality and are significantly more efficient than traditional loop-based implementations. Bellmania includes a high-level language for specifying dynamic programming algorithms and a calculus that facilitates gradual transformation of these specifications into efficient implementations.
These transformations formalize the divide-and-conquer technique; a visualization interface helps users to interactively guide the process,
while an SMT-based back-end verifies each step
and takes care of low-level reasoning required for parallelism.

We have used the system to generate provably correct implementations of several algorithms, including some important algorithms from computational biology, and show that the performance is comparable to that of the best manually optimized code.
\end{abstract}

\category{D.1.2}{Programming Techniques}{Automatic Programming}
 
% general terms are not compulsory anymore, 
% you may leave them out
%\terms
%term1, term2
\keywords{Divide-and-conquer, Dynamic Programming, Program Synthesis, Verification, Refinement Types, SMT}
%\keywords
%keyword1, keyword2

\begin{wrapfigure}{R}{3cm}
\vspace{-1.35in}\includegraphics[scale=0.1]{img/aec-badge}
\end{wrapfigure}

\section{Introduction}
\label{intro}

\newcommand\ForLine[1]{\State\algorithmicfor~{#1} \algorithmicdo~}
\newcommand\Head[1]{\Comment{ {\it #1} ~~}}

Software synthesis aims to close the gap between descriptions of software components, such as algorithms and systems, and their implementations as computer programs. Dynamic Programming (DP) algorithms offer a prominent example of how large this gap can be. For instance, consider  
\cbstart\diffnote{some short explanation, and link}%
\Cref{intro:naive}, which correspond to the well known DP algorithm to compute the optimal parenthesization for a chain of matrix multiplications.\footnote{\fontsize{6}{6}\selectfont\url{https://en.wikipedia.org/wiki/Matrix_chain_multiplication}}
The first loop fills the diagonal of an array with some initial values,
and the second loop computes off-diagonal elements by reading existing elements.
\cbend

The algorithm computes an $n\times n$ region of a DP table $G$ via a standard row-major order. This algorithm is simple, but a direct C implementation of it turns out to be about
$10\times$ slower than the best manually optimized implementation \cbstart(as we will see in \Cref{evaluation:experimental})\cbend. 
One reason for this poor performance is the high rate of cache misses incurred by reading the ranges $G_{ik}$ and $G_{kj}$, for $i<k<j$, repeatedly on every iteration of the loops over $i$ and $j$. Memory reads dominate the running time of this algorithm, so high speedups can be gained by localizing memory access.

\begin{algorithm}[b]
\renewcommand\arraystretch{1.3}
\begin{algorithmic}
  \ForLine{$i=1..n$}  $G_{i(i+1)} := x_i$    \Head{Initialize}
  \For{$i=(n-2)..0$}          \Head{Compute}
    \For{$j=(i+2)..n$}
      \State $G_{ij} := \underset{i<k<j}\min ~ G_{ik} + G_{kj} + w_{ikj}$
    \EndFor
  \EndFor
\end{algorithmic}
\caption{\label{intro:naive}
   A na\"ive loop implementation}
\end{algorithm}


The state-of-the-art implementation uses a \newterm{divide and conquer} approach both to improve memory performance and to increase the asymptotic degree of parallelism~\cite{IPDPS15/Tithi}. An excerpt of the pseudo-code for such an implementation is shown in \Cref{intro:divide-and-conquer}. In the optimized version, the programmer has to determine a number of low-level details, including the correct order of calls---some of which can be run in parallel---as well as some rather involved index arithmetic. When implemented in C++, the full version is, in fact, more than ten times longer than the na\"ive one and considerably more complicated. It is also much more difficult to verify, since the programmer would have to provide invariants and contracts for a much larger set of loops and possibly recursive functions. Parallelism only adds to the complexity of this task.



\begin{algorithm}[b!]
\fontsize{9}{10}   % being sneaky!
\begin{algorithmic}
  \State A[$1..n$], ~where: \hfill{\normalsize (\textit{snippet})}
  \Procedure{A{\larger{[}}$s..e${\larger{]}}}{}
	\If{$e-s < b$}
	  \For{$i=e..s$}
	    \For{$j=\max\{s,i\}..e$}
	      \State $G_{ij} := \underset{i<k<j}\min ~ G_{ik} + G_{kj} + w_{ikj}$
	    \EndFor
	  \EndFor
    \vspace{-5pt}
	\Else
	  \State A$\big[ s..\floor{\frac{s+e}{2}} \big]$
	  \State A$\big[ \floor{\frac{s+e}{2}}{+}1..e \big]$
	  \State B$\big[ s..\floor{\frac{s+e}{2}} ~,~ \floor{\frac{s+e}{2}}{+}1..e \big]$
	\EndIf
  \EndProcedure
  \Procedure{B{\larger{[}}$s_0..e_0 ~,~ s_1..e_1${\larger{]}}}{}
	\If{$e-s < b$} \ldots
	\Else
	  \State B$\big[ \floor{\frac{s_0+e_0}{2}}{+}1..e_0 ~,~ s_1..\floor{\frac{s_1+e_1}{2}} \big]$
	  \State C$\big[ s_0..\floor{\frac{s_0+e_0}{2}} ~,~ s_1..\floor{\frac{s_1+e_1}{2}} ~,~ \floor{\frac{s_1+e_1}{2}}{+}1..e_1 \big]$
	  \vspace{-5pt}
	  \State $\vdotswithin{C}$
	\EndIf
  \EndProcedure
  \Procedure{C{\larger{[}}$s_0..e_0 ~,~ s_1..e_1 ~,~ s_2..e_2${\larger{]}}}{}
    \vspace{-5pt}
    \State $\vdots$
  \EndProcedure
\end{algorithmic}
\caption{\label{intro:divide-and-conquer}
   An optimized divide-and-conquer version}
\end{algorithm}




In this paper, we present a new system called Bellmania\footnote{Named so as a tribute to Richard Bellman.}, which allows an expert to interactively generate parallel divide-and-conquer implementations of dynamic programming algorithms, which are provably correct relative to a high-level specification of the code in \Cref{intro:naive}. We show that the resulting implementations are \cbstart\diffnote{use ranges}1.4--46$\times$ faster than the original versions and within 2--60\%\cbend{} of a high-performance hand crafted implementation.
Their structure will also make them \newterm{cache oblivious}~\cite{FOCS99/Frigo} and \newterm{cache adaptive}~\cite{SODA14/Bender}, just like the hand crafted implementations are.


Bellmania is a deductive synthesis system in the tradition of systems like KIDS~\cite{TSE90/Smith}, and more recently Fiat~\cite{POPL15/Delaware}. 
These systems derive an implementation from a specification in a correct-by-construction manner by applying a sequence of deductive reasoning steps. Thanks to the correct-by-construction approach, the potentially complex final artifact does not have to be independently verified, making the approach ideal for potentially complex implementations like the ones we target. 

Traditionally, the major shortcoming of deductive synthesis has been the effort required of the user in order to guide the derivation process towards a correct solution. In this work, we reduce this effort using three core techniques. 
First, we show that a small number of domain specific tactics, combined with a new notation to jointly describe the specification and the implementation, can enable the derivation to be described succinctly at a high-level of abstraction. Secondly, we show the value of incorporating an SMT solver into the derivation process; in particular, we can afford to use tactics that are only correct when some side conditions hold, where these conditions are potentially complex logical assertions, without having to burden the user with proving those premises. Finally, we show that by incorporating \emph{solver-based inductive synthesis},
which generalizes from concrete values and execution traces, into the derivation process, we can automate away many low-level decisions, allowing for shorter, simpler derivations. We use the term \emph{solver-aided tactics} to refer to this combination of solver-based inductive synthesis and solver-checked premises within a tactic.

Overall, the paper makes the following contributions. 
\begin{itemize}
\item A new formalism used to describe a wide class of dynamic programming algorithms, capable of bridging the gap between the high-level specification and the divide-and-conquer implementation of them.

\item An interesting application of refinement types for tracking dependencies between sub-computations, making it possible to automatically generate parallel implementations from it.

\item The idea of using \emph{solver-aided tactics}, demonstrating their applicability and utility in the derivation of divide-and-conquer dynamic programming implementations. 

\item A suite of solver-aided tactics for dynamic programming and an overview of the proofs of their soundness, assuming only the soundness of the underlying SMT solver. 

\item A description of Bellmania, the first system capable of generating provably correct implementations of di\-vide-and-conquer dynamic programming. Our evaluation shows that the code it generates is comparable to manually tuned implementations written by experts in terms of performance.
\end{itemize}

Dynamic Programming is central to many important domains ranging from logistics to computational biology --- e.g.,~a recent textbook \cite{DurbinEdKr98} lists 11 applications of DP in bioinformatics just in its introductory chapter, with many more in chapters that follow. Increasing performance and reliability of DP implementations can therefore have significant impact. More generally, we believe that this work serves as an important test case for a new approach to combining inductive and deductive synthesis which could have an impact beyond this domain. 



%==============================================



%In this paper, we present a new approach to interactive deductive synthesis based on \emph{solver-aided tactics} that preserves the benefits of deductive synthesis techniques but reduces the burden on the user by relying on two important innovations: (a) the use of inductive synthesis to discover important details of the low-level steps needed to achieve a transformation, (b) the use of a type system based on predicate abstraction (liquid types) that associates semantic information with program terms, enabling automated verification of the validity of a transformation. Both of these innovations rely on extensive use of SMT solvers to discharge complex proof obligations that would otherwise have to be discharged interactively with significant manual effort. 


\begin{comment}


Our work on solver-aided tactics builds on prior work on the StreamBit project~\cite{PLDI05/Solar-Lezama}, which
introduced the idea of transformation rules with missing details that can be inferred by a symbolic search procedure, as
well as the pioneering work on the Leon synthesizer\coa{cite?}, which has explored the integration of deductive techniques to improve the
scalability of inductive synthesis. However, our approach is unique in the way it leverages inductive synthesis and
typing in the context of deductive synthesis: 
(a) the solver can use inductive synthesis to search for the detailed parameters required for a transformation,
(b) The solver can prove validity of side conditions that ensure the soundness of each individual transformation, 
(c) the tactics can leverage information from \newterm{logically qualified types} in the program in
    order to guide the transformation. 
Being able to rely on the solver to check the validity of transformations gives greater flexibility when designing such transformations. 
It means that we do not have to be conservative when stating preconditions.
In addition, the reliance on the synthesizer to fill in parameters and proof details means that the user has to type less and has fewer chances to make mistakes.

Overall, we make the following contributions.
\begin{itemize}
\item We introduce \emph{solver-aided tactics}, a form of rewrite rules with proof obligations that can be translated to first-order logic,
  as a way to raise the level of abstraction of deductive synthesis.
\item We develop a small library of these formal tactics that can be used to 
  systematically transform a class of problem specifications,
  expressed as recurrences in a simple functional language,
  into equivalent divide-and-conquer programs that admit cache-oblivious parallel
  implementations.
\item We prove that these tactics are semantics-preserving, assuming some side conditions are met
  at the point when the tactic is applied.
\item We show that the side conditions can be effectively translated into first-order closed
  formulas, and verified automatically by SMT solvers.
\item We demonstrate the first system capable of generating provably correct implementations of divide-and-conquer implementations from a high-level description of the algorithm. 
\item We measure the performance of automatically generated code and show that it is comparable to manually tuned reference implementations written by experts.
\end{itemize}


\end{comment}


%We believe the approach has the potential to be generally applicable to a variety of synthesis problems, but in this paper, we focus on a particular domain of \emph{divide-and-conquer dynamic programming} algorithms. Specifically, we have developed a system called \emph{Bellmania} that uses solver-aided tactics specialized for this domain to help an algorithm designer derive divide-and-conquer dynamic programming algorithms from a high-level specification. 


\section{Overview}
\label{overview}

Most readers are likely familiar with the Dynamic Programming (DP) technique of Richard Bellman~\cite{03/Bellman:DP} to construct an optimal solution to a problem by combining together optimal solutions to many overlapping sub-problems. The key to DP is to exploit the overlap and reuse computed values to explore exponential-sized solution spaces in polynomial time. Dynamic programs are usually described through recurrence relations that specify how to decompose sub-problems, and is typically implemented using a DP table where each cell holds the computed solution for one of these sub-problems. The table can be filled by visiting each cell once in some predetermined order, but recent research has shown that it is possible to achieve order-of-magnitude performance improvements over this standard implementation approach by developing \emph{divide-and-conquer}  implementation strategies that recursively
partition the space of subproblems into smaller subspaces~\cite{IPDPS15/Tithi,SODA14/Bender,SODA06/Chowdhury,SPAA08/Chowdhury,TOCS10/Chowdhury,TCBB10/Chowdhury}. 

Before delving into how Bellmania supports the process of generating such an implementation, it is useful to understand how a traditional iterative implementation works. For this, we will use the 
optimal parenthesization algorithm from the introduction (\Cref{intro:naive}). The problem description is as follows: given a sequence of factors $a_0\!\cdots a_{n-1}$, 
the goal is to discover a minimal-cost placement of parentheses in the product expression
assuming that multiplication is associative but not commutative. The cost
of reading $a_i$ is given by $x_i$, and that the cost for multiplying
$\Pi(a_{i..(k-1)})$ by $\Pi(a_{k..(j-1)})$ is given by $w_{ikj}$. The specification of the algorithm is shown in \Cref{overview:paren spec}; 
the values $x_i$ and  $w_{ikj}$ are inputs to the algorithm, and the output is a table $G$, where each element $G_{ij}$ is the lowest cost for parenthesizing $a_{i..(j-1)}$, with $G_{0n}$ being the overall optimum.

\begin{figure}[b]
\begin{tabular}{@{\hspace{-4pt}}r@{~~~}p{5cm}@{\hspace{-4pt}}}
\begin{tikzpicture}[x=4.1mm,y=4.1mm,baseline=(center), remember picture]
  \coordinate(center) at (3,3);
  \def\n{7}  % actually n+1 but let's not be nitpicky
  \draw[black!25!white] (\n,1) |- (0,0) |- (1,\n);
  \draw (\n,1) -- (\n,\n) -- (1,\n);
  \foreach \i [evaluate=\i] in {2-1,...-1,\n-1} {
    \draw[black!25!white] (0,\i) -- (\n-\i,\i); 
    \draw[black!25!white] (\i,0) -- (\i,\n-\i); 
    \draw (\n,\i) -- (\n-\i,\i); 
    \draw (\i,\n) -- (\i,\n-\i);
  }
  %\draw[step=1] (0,0) grid (6,6);
  \draw[ultra thick] (5,5) rectangle +(1,1);
  \node[circle,fill=BrickRed,inner sep=0,minimum size=1mm](Gij) at (5.5,5.5) {};
  \fill[black,opacity=0.1] (1,\n) rectangle ++(1,-1) rectangle ++(1,-1) rectangle ++(1,-1) rectangle ++(1,-1) rectangle ++(1,-1) rectangle ++(1,-1);
  \fill[blue,opacity=0.2] (2,5) rectangle (5,6);
  \fill[blue,opacity=0.2] (5,2) rectangle (6,5);
  \node[anchor=south east](G) at (0,\n) {\small$G$};
  \draw[->] (G.east) -- +(1.5,0) node[anchor=west] {\small $j$};
  \draw[->] (G.south) -- +(0,-1.5) node[anchor=north] {\small $i$};
\end{tikzpicture}
&
\vspace{-1.5cm}
\small
$
\begin{array}{@{}l@{}}
	\tikz[overlay, remember picture]{\draw[BrickRed] (-.5mm,0) -- (Gij);}
	G_{ij} ~=~ \\
	~
	\begin{cases}
		~~x_i                        & i{+}1=j \\
	    \underset{i<k<j}\min ~ G_{ik} {+} G_{kj} {+} w_{ikj} & i{+}1<j
	\end{cases}
\end{array}
$

\vspace{6mm}
% legend
\cbstart
\begin{tikzpicture}
  \node(base)[draw,fill=black!10!white,label=right:{\fontsize{8}{8}\selectfont base case}] {};
  \node(dep)[right=1.5cm of base,draw,fill=blue!20!white,label=right:{\fontsize{8}{8}\selectfont dependencies}] {};
  \node(emp)[right=1.8cm of dep,draw=black!25!white,label=right:{\fontsize{8}{8}\selectfont empty}] {};
\end{tikzpicture}
\cbend
\end{tabular}
\caption{Recurrence equation and cell-level dependencies.}
\label{overview:paren spec}
\end{figure}

\begin{paragraph}{Iterative Algorithm.}
Using the standard dynamic programming method, anyone who has  read \cite{09/CLRS} would compute this recurrence
with an iterative program by understanding the dependency pattern:
to compute the $\min_{i<k<j}(\cdots)$ expression in \Cref{overview:paren spec}
the algorithm needs to enumerate $k$ and gather information from all cells below and to the left of $G_{ij}$.
In particular, each value $G_{ij}$ is computed from other values $G_{i'j'}$ with higher
row indexes $i'>i$ and lower column indexes $j'<j$. 
Therefore, considering $G$ as a two-dimensional array, it can be filled in a single sweep from left to right and from bottom
to top, as done in \Cref{intro:naive}.
\end{paragraph}




\newcommand\qbox[1]{\fbox{\rm\scriptsize#1}}
\newcommand\tinyqbox[1]{\hspace{.5pt}\tikz \node[draw,inner sep=1.5pt] {$\scriptscriptstyle #1$};}

\newcommand\plusoneocd{\raisebox{.5pt}{$\scriptstyle+1$}}

\algrenewtext{Procedure}{\hspace{-3mm}{\bf procedure}~}  % hack to make proc header slighly less indented

\begin{paragraph}{Divide-and-Conquer Algorithm.}
To illustrate the main concepts underlying Bellmania and the key ideas in deriving divide-and-conquer implementations, 
we will walk through the first
few steps that an algorithms expert --- whom we will call Richard --- would follow using Bellmania to generate a provably correct divide-and-conquer implementation of the optimal parenthesization algorithm.

In the Bellmania development model, Richard will start with the specification from \Cref{overview:paren spec}, and progressively manipulate it to get the specification in a form that reflects the recursive structure of the divide-and-conquer implementation. At any step in the transformation, Bellmania can generate code from the partially transformed specification. Code generated from the initial specification will yield an implementation like \Cref{intro:naive}, whereas generating code from the fully transformed specification will yield the divide-and-conquer implementation that we want.
In the rest of the text, we will use the term \newterm{program} to refer to any of the partially transformed specifications.




\begin{figure}
\centering
\begin{tabular}{c@{\hspace{.5in}}c}
\begin{tikzpicture}[baseline=(n/2), q/.style={font=\relsize{1.3}}, pale/.style={color=black!25!white}]
  % this code draws a grid with the bottom-left paler
  %\draw (0,2) -- (2,2) -- (2,0);
  %\foreach \i in {0,1} {
  %  \draw[pale] (\i,0) -- (\i, 1-\i); \draw (\i,2) -- (\i, 1-\i);
  %  \draw[pale] (0,\i) -- (1-\i, \i); \draw (2,\i) -- (1-\i, \i);
  %}
  % but this might be confusing. let's just paint
  % a regular 2x2 grid
  \draw (0,0) grid (2,2);
  \node[q] at (.5,1.5) {1};   \node[q] at (1.5,1.5) {2};
  \node[q] at (.5, .5) {3};   \node[q] at (1.5, .5) {4};
  \node(O)[above left] at (0,2) {$0$};
  \node(m/2)[above] at (1,2) {$\frac{n}{2}$};
  \node(m)[above] at (2,2) {$n$};
  \node(n/2)[left] at (0,1) {$\frac{n}{2}$};
  \node(n)[left] at (0,0) {$n$};
  \node(J0)[above] at (.5,2.5) {$J_0$};
  \node(J1)[above] at (1.5,2.5) {$J_1$};
  \node(I0)[left] at (-.5,1.5) {$J_0$};
  \node(I1)[left] at (-.5,.5) {$J_1$};
  %\coordinate(0) at (0,0);
  %\coordinate(sw) at (0,0);
  %\coordinate(ne) at (2,2);
  %\draw (J0.north -| sw) -- node[above] {$J$} ++(ne |- 0);
  %\draw (I0.west |- sw) -- node[left] {$I$} ++(ne -| 0);
  \draw (O.north east) -- (O.north east -| m/2.north west);
  \draw (O.north east -| m.110) -- (O.north east -| m/2.north east);
  \draw (O.south west) -- (O.south west |- n/2.north west);
  \draw (O.south west |- n.160) -- (O.south west |- n/2.220);
\end{tikzpicture}
& 
$\begin{array}{l}
  \qbox2 \mbox{\textit{ depends on }} \qbox1 \\[2pt]
  \qbox2 \mbox{\textit{ depends on }} \qbox4 \\ 
  \\
  (\,\qbox3\mbox{\textit{ is empty}})
\end{array}$
\end{tabular}
\vspace{5pt}
\caption{\label{overview:quadrants}
  Dividing a two-dimensional array into quadrants; the dependencies for the case of the running example are shown on the right.}
\end{figure}


\Cref{overview:slice-stratify-synth} provides a visual description of the initial stages of the transformation process. The figure includes block diagrams illustrating how the program at a given stage in the transformation will compute its output table from its input. For example, the first row corresponds to the program before any transformations take place, i.e.~the initial specification. At this stage, the program involves a single loop nest that reads from the entire array and writes to the entire array; solid arrows in the diagram denote data dependency. The transformation from one stage to the next, the dashed arrows, is performed by the application of  \newterm{tactics} that represent a high-level refinement concept. 

As a first step in the transformation, Richard would like to partition the two-di\-men\-sio\-nal array $G$ into quadrants, as illustrated in \Cref{overview:quadrants}. In Bellmania, the partition is accomplished by applying the {\sf Slice} tactic, illustrated graphically at the top of \Cref{overview:slice-stratify-synth}. In order to escape the need to reason about concrete array indices, Bellmania provides an abstract view where
the partitions are labeled $J_0,J_1$. The effect of {\sf Slice} is shown in text in \Cref{overview:logical-slice-stratify}$(a)$ --- the figure trades accuracy for succinctness by omitting the base case of the recurrence for now. Due to the structure of the problem, namely $i\!<\!j$, 
the bottom-left quadrant (\qbox3 in \Cref{overview:quadrants}) is empty.
Slicing therefore produces only three partitions.


The computation of \qbox1
(the top-left quadrant)
does not depend on any of the other computations, so Richard applies the
{\sf Stratify} tactic, which separates an independent computation step as a separate loop. This is equivalent to rewriting the specification as in \Cref{overview:logical-slice-stratify}({\it b}):
the first computation is given a special name $G^{\tinyqbox1}$, then the following
computations read data either from $G^{\tinyqbox1}$ (when the indices are in \qbox1)
or from $G$ (otherwise), which is denoted by $G^{\tinyqbox1}\!/G$. The ``$/\,$'' operator
is part of the Bellmania language and will be defined formally in \Cref{lang}.
Bellmania checks the data dependencies and verifies that the transformation
is sound.

Repeating {\sf Stratify} results in a three-step computation,
as seen in \Cref{overview:slice-stratify-synth}$(c)$, from which Richard can obtain the program in \Cref{overview:breakdown}.
This already gives some performance gain, since computations of \qbox1 and \qbox4 can now run in parallel.
Bellmania is capable of sound reasoning about parallelism, using knowledge encoded
via types of sub-terms, showing that two threads are race-free when they work on different regions
of the table.
This is handled automatically by the code generator.
\end{paragraph}


\begin{figure}
\vspace{-.5em}
\[\renewcommand\arraystretch{1.3}
  \begin{array}{@{}l@{}}
    \textsf{Slice} \quad i,j:\langle J_0\times J_0 \,|\, J_0\times J_1 \,|\, J_1\times J_1\rangle     \hfill (a)\\
    ~~\forall i,j\in\qbox1.~ G_{ij} = \underset{i<k<j}\min ~ G_{ik} + G_{kj} + w_{ikj} \\
    ~~\forall i,j\in\qbox2.~ G_{ij} = \underset{i<k<j}\min ~ G_{ik} + G_{kj} + w_{ikj} \\
    ~~\forall i,j\in\qbox4.~ G_{ij} = \underset{i<k<j}\min ~ G_{ik} + G_{kj} + w_{ikj} \\
    %
    \textsf{Stratify} ~ \qbox1 \hfill (b)\\
    %
    ~~\forall i,j\in\qbox1.~ G^{\tinyqbox1}_{ij} = \underset{i<k<j}\min ~ G^{\tinyqbox1}_{ik} + G^{\tinyqbox1}_{kj} + w_{ikj} \\
    ~~\forall i,j\in\qbox2.~ G_{ij} = \underset{i<k<j}\min ~ (G^{\tinyqbox1}\!/G)_{ik} + (G^{\tinyqbox1}\!/G)_{kj} + w_{ikj} ~~ \\
    ~~\forall i,j\in\qbox4.~ G_{ij} = \underset{i<k<j}\min ~ (G^{\tinyqbox1}\!/G)_{ik} + (G^{\tinyqbox1}\!/G)_{kj} + w_{ikj} \\
  \end{array}
\]
\vspace{-.75em}
\caption{\label{overview:logical-slice-stratify}
  The first two steps in the development, represented as logical specifications.}
\end{figure}


\newcommand\steparrowwidth{3mm}
\newcommand\steparrow{\includegraphics[width=\steparrowwidth]{img/arrow}}

\begin{algorithm}[b]
\renewcommand\arraystretch{1.3}
\begin{algorithmic}
\Procedure{A{\larger{[}}$G${\larger{]}}}{}\EndProcedure
  \For{$i=(n{-}2)..0 \cap J_0$}    \Head{Compute \qbox1}
    \For{$j=(i{+}2)..n \cap J_0$} 
      \State $G_{ij} :=
          \underset{i<k<j}\min ~ G_{ik} + G_{kj} + w_{ikj}$
    \EndFor
  \EndFor
  \For{$i=(n-2)..0 \cap J_1$}    \Head{Compute \qbox4}
    \For{$j=(i{+}2)..n \cap J_1$} 
      \State $G_{ij} :=
          \underset{i<k<j}\min ~ G_{ik} + G_{kj} + w_{ikj}$
    \EndFor
  \EndFor
  \For{$i=(n{-}2)..0 \cap J_0$}    \Head{Compute \qbox2}
    \For{$j=(i{+}2)..n \cap J_1$}
      \State $G_{ij} :=
          \underset{i<k<j}\min ~ G_{ik} + G_{kj} + w_{ikj}$
    \EndFor
  \EndFor
\end{algorithmic}
\caption[t]{\label{overview:breakdown}
   \cbstart Parenthesis\cbend{} --- Sliced and Stratified}
\end{algorithm}


\newbox\primebox
\setbox\primebox\hbox{$'$}
\newbox\doubleprimebox
\setbox\doubleprimebox\hbox{$''$}

\newcommand\primeocd[1]{\hspace{\wd\primebox}#1\usebox\primebox}
\newcommand\doubleprimeocd[1]{\hspace{\wd\doubleprimebox}#1\usebox\doubleprimebox}

\begin{figure}
\centering
\input{gfx/overview-slice-stratify-synth}
\caption[caption]{\label{overview:slice-stratify-synth}
  Overview of tactic semantics in Bellmania.}
\end{figure}


\medskip
At this point, Richard notices that $G^{\tinyqbox1}$ is just a smaller version of
the entire array $G$; he invokes another tactic called {\sf Synth}, which automatically
synthesizes a recursive call A\big[$G^{\tinyqbox1}$\big]
(presented using abstract index ranges as $A^{J_0}$).

Similarly, $G^{\tinyqbox4}$ is also a smaller version of $G$, this time with
indices from $J_1$. {\sf Synth} figures it out automatically as well,
synthesizing a call A\big[$G^{\tinyqbox4}$\big]. The remaining part,
$G^{\tinyqbox2}$, is essentially different, so Richard gives it a new name, ``B'',
which becomes a new synthesis task.%
\footnote{Richard's choice of names is consistent with the literature.}

\begin{algorithm}[b]
\renewcommand\arraystretch{1.3}
\begin{algorithmic}
\Procedure{A{\larger{[}}$G${\larger{]}}}{}\EndProcedure
  \If{$G$ {\it is very small}} {\it run iterative version}
  \Else
  \State A\big[$G^{\tinyqbox1}$\big] \Head{Compute \qbox1}
  \State A\big[$G^{\tinyqbox4}$\big] \Head{Compute \qbox4}
  \State B\big[$G^{\tinyqbox1}, G^{\tinyqbox4}, G^{\tinyqbox2}$\big] \Head{Compute \qbox2}
  \EndIf
\end{algorithmic}
\caption{\label{overview:recursive-A}
   Parenthesis --- Recursive Version}
\end{algorithm}


Applying the same strategy will eventually lead Richard to further
break down and transform the computation of B into multiple recursive
sub-computations, %\rch{B makes eight recursive function calls to update the four quadrants. So it reduces to eight recursive sub-computations, right?}, 
further improving the locality of the resulting algorithm until
a true divide-and-conquer solution is obtained.
Bellmania generates code for
all procedures encountered throughout the development. In this case, three recursive procedures
are generated. The base case of the recursion is when the region
becomes small enough to just run the loop version.\footnote{The optimal  base case size of the region can be found by auto-tuning (taken as an input in the current version of the compiler)}

As is well illustrated by the example, this line of reasoning can get quite complicated for most dynamic programming algorithms, 
and producing a correct divide-and-conquer algorithm for a given dynamic programming problem is considered quite difficult even by the researchers who originally pioneered the technique. 
Fortunately, Bellmania is able to mechanize most of the technical details,
allowing Richard and other algorithm designers to focus on their area of expertise,
try different strategies,
and eventually produce a \emph{certified} implementation of the algorithm.

Overall, it took Richard only 4 steps to construct \Cref{overview:recursive-A},
and a total of 30 steps to construct all three phases of the Parenthesis algorithm,
comprising an implementation that is 46$\times$ faster than a parallel implementation of
\Cref{intro:naive} using a state-of-the-art parallelizing compiler.
The user is greatly assisted by tactics like {\sf Synth}, which carries out the monotonic
and error-prone task of choosing the right parameters for each recursive call; also,
mistakes are identified early in the development thanks to automatic verification,
saving hours of debugging later on.
The resulting code is much easier to maintain, because the artifact is not just
the optimized C++ code, but also the Bellmania specification and the script
that encodes the optimization strategy.

Once a divide-and-conquer algorithm is found, generating an optimal implementation still requires some additional work, such as finding the right point at which to switch to an iterative algorithm to leverage SIMD parallelism as well as low-level tuning and compiler optimization;
these steps are performed by more traditional compiler optimization techniques
as discussed in \Cref{codegen}.


\subsection{System Design}

\begin{figure}
  \centering
  
  \begin{tikzpicture}[
      h/.style={minimum height=1.75em}, 
      tallh/.style={minimum height=3.1em}, 
      doubleh/.style={minimum height=3.5em+2mm}, 
      fullw/.style={minimum width=5cm}, 
      halfw/.style={minimum width=2.4cm},
      75w/.style={minimum width=1.15cm},
      every node/.style={align=center}]
    \node[draw, h, fullw](tae) {Tactic Application Engine};
    \node[draw, h, fullw, above=2mm of tae](dp) {\fontsize{9.5}{11}\selectfont Divide-and-Conquer Tactic Library};
    \node[draw, tallh, halfw, below=2mm of tae.south west, anchor=north west](verif) {Verifier \\ {\small (SMT)}};
    \node[draw, tallh, halfw, below=2mm of tae.south east, anchor=north east](comp) {Compiler \\ {\small (C++, cilk)}};
    \draw (dp.south) -- (tae.north);
    \draw (tae.south -| verif.north) -- (verif.north);
    \draw (tae.south -| comp.north) -- (comp.north);
    
    \node[draw, doubleh, left=4mm of dp.north west, anchor=north east](ui) {Frontend \\ UI};
    \draw (dp.west) -- (dp.west -| ui.east);
    \draw (tae.west) -- (tae.west -| ui.east);
    
    \node[draw, tallh, right=4mm of comp](exe) {Exec. \\ Code};
    \draw[double] (comp) -- (exe);
  \end{tikzpicture}
\caption{\label{overview:design}
  Overall design of Bellmania.}
\end{figure}

The design of the Bellmania system (\Cref{overview:design}) contains a generic
core --- called the \newterm{Tactic Application Engine} (TAE) ---
on top of which a library of \newterm{tactics} specific to dynamic programming
is built. While it is possible to extend the library, it already contains enough tactics to successfully develop algorithms for a
family of problems, so that the user of the system only needs to apply existing tactics by issuing commands
through a GUI and watching the program evolve. The TAE has a back-end that
verifies conjectures, and is in charge of making sure that
tactic applications represent valid rewritings of the program. Finally, the
programs created this way are transferred to a compilation back-end, where
some automatic static analysis is applied and then executable (C++) code is emitted.

\cbstart\diffnote{made the point clear regarding verification vs. certification, and what is missing}%
The trusted core is small, comprising of:
\begin{enumerate*}[label=(\textit{\alph*})]
 \item a type checker (see \Cref{lang:types}), 
 \item a term substitution procedure (see \Cref{tactics}),
 \item a formula simplifier (see \Cref{automated:simplification}),
 \item the SMT solver, and
 \item the compiler (see \Cref{codegen}).
\end{enumerate*}
In its current version, Bellmania does not emit machine-checkable
proofs; to the extent that SMT solvers can emit such proofs,
these could be adapted to a proof of refinement (in the sense of \cite{POPL15/Delaware}).
The compiler has to be formally verified separately.
These two tasks require considerable engineering effort and are left for the future.
\cbend

\begin{figure}
\vspace{-.5em}% why is there extra space above lstlisting?
\begin{lstlisting}[language=bellmania]
Slice (find (\theta |-> ?)) (? <|J_0*J_0, J_0*J_1, J_1*J_1|>)
Stratify "/" (fixee [A]) [A] \psi
Stratify "/" (fixee [A]) [A] \psi
[A] [B] [C] |-> SynthAuto . ... \psi
\end{lstlisting}
\caption{\label{overview:script}
  Bellmania script used to generate \Cref{overview:recursive-A}.}
\end{figure}

An example for the concrete syntax is shown in \Cref{overview:script}.
A full listing of the scripts for our running examples,
as well as screenshots of the UI, can be found
in \Cref{annex:example:paren}.

\subsection{Interaction Model}

\begin{figure}[b!]
\begin{tikzpicture}[>=latex,
     blk/.style={draw,rectangle,align=center,inner sep=5pt,minimum height=8mm}]
  \node(spec)[blk]{Specification};
  \node(prog)[blk,below right=8mm of spec]{Program Term};
  \node(deriv)[blk,above=of prog]{Derived\\Definition};
  \node(impl)[blk,above right=8mm of prog]{Implementation};
  \draw[->,rounded corners=5mm] (spec) |- node[below]{\textit{typecheck}} (prog);
  \draw[->] (prog) edge[out=112,in=-105] (deriv);
  \draw[->] (deriv) edge[out=-73,in=70] (prog);
  \draw[->,rounded corners=5mm] (prog) -| node[below]{\textit{compile}} (impl);
  % Loop over (prog)
  \draw[->,rounded corners=3mm] (prog.-10) -| ++(.3,-.6) -| ($(prog.190) + (-.3,0)$) -- (prog.190);
  \node[below=6mm of prog]{\textit{apply tactic}};
\end{tikzpicture}
\caption{\label{overview:flow}
  Interaction workflow in Bellmania.}
\end{figure}

\cbstart\diffnote{added subsection (incl. figure)}%
The intended use pattern for Bellmania is by repeatedly issuing
commands to apply tactics inside a REPL that keeps showing to the
user the resulting \newterm{program term} (\Cref{overview:flow}).
To begin the interaction, the user types in a \newterm{specification} written
in the Bellmania language (as defined in \Cref{lang}), which is typechecked
by the system, producing a term.
This term then becomes the focus of the development, and further transformations
apply to it; the user issues \newterm{apply tactic} commands, one at a time,
using a syntax similar to \Cref{overview:script}.
Tactic applications are also typechecked, as well as the resulting program,
since type information can be refined by including the context (see \Cref{lang:types}).
During the development, sub-computations may be discovered that require
further drilling down (such as ``B'' in the last example).
These result in \newterm{derived} terms that take the focus until they, too,
have been fully transformed.
When the program and all its sub-computations are fully developed and
expressed as a divide-and-conquer algorithm, the user invokes the compiler
back-end that emits C++ code.

In the following sections, we describe the programming language and
formally define the tactics that were used in the example above.
We then show how to formalize the same intuition as we had there,
using this new instrument.
\cbend

\section{A Unified Language}
\label{lang}

\newcommand\semp[1]{[\![{#1}]\!]}
\newcommand\fix{\operatorname{fix}}

For the purpose of manipulating programs and reasoning about them,
we first set up a small language that can easily express our programs (that is,
partially transformed specifications), and can also be translated to executable
code in a straightforward fashion. We choose a functional
setting, with scalar types for array indices ($i, j, k : J$ in the example)
and values stored in the arrays, typically real numbers ($G_{ij} : \R$ in the example).
Arrays are encoded as functions, or arrow types, mapping indices
to values (e.g. $G : J^2\to\R$). We want to have a notion of
array elements that are uninitialized, so we assume every scalar type
contains the special value $\bot$, corresponding to an undefined value.

Our language should be able to express operations such as \textsf{Slice}
from the previous section. To achieve this, we introduce 
\newterm{predicate abstraction} over the index types, by extending the
type system with subtyping and \newterm{refinement types}. An index
type $J$ can be partitioned into subtypes $\langle J_0 \,|\, J_1\rangle$
(two or more), meaning that we define predicates $\widehat{J}_0,\widehat{J}_1 : J\to\B$ (where $\B$ is the Boolean type)
and the refinement types $J_0=\{v\,{:}\,J \,|\, \widehat{J}_0(v)\}$
and $J_1=\{v\,{:}\,J \,|\, \widehat{J}_1(v)\}$, which become subtypes of $J$.
Usually, we can omit the ``hat'' and use $J_0$ as a synonym for $\widehat{J}_0$ 
when it is clear from the context that it designates a type. 
We would then provide types for sub-arrays, e.g. 
$G^{\raisebox{-2pt}{\tinyqbox{2}}} : J_0\times J_1 \to \R$.

To refer to a region of a given array, we define a \newterm{guard}
operator, which is parameterized by subtypes of the indices,
for example $G^{\raisebox{-2pt}{\tinyqbox{2}}} = [G]_{J_0\times J_1}$.
To combine computations on different parts of the array, we use a lifting
of the ternary condition operator (known as {\tt ?:}) to functions,
where the condition is implied by the operands; e.g.
(\Cref{overview:quadrants})
\[
G=G^{\tinyqbox1} \big/ G^{\tinyqbox2} \big/ G^{\tinyqbox4} = 
\lambda ij.
\begin{cases}
  G^{\tinyqbox1}_{ij} & i,j\in J_0 \\
  G^{\tinyqbox2}_{ij} & i\in J_0 \land j\in J_1 \\
  G^{\tinyqbox4}_{ij} & i,j\in J_1 \\
\end{cases}
\]

\bigskip
\begin{paragraph}{Formal set-up}The Bellmania language is based on the polymorphic
$\lambda$-calculus, that is, simply typed $\lambda$-calculus with universally quantified type variables 
(also known as \newterm{System F}).
The following subsections contain formal definitions for language constructs.
\end{paragraph}

We write abstraction terms as $(v:\T)\mapsto e$, where $\T$ is the type of the argument $v$ and $e$ is
the body, instead of the traditional notation $\lambda(v:\T).\,e$, mainly due to aesthetic reasons
but also because we hope this will look more familiar to intended users.
Curried functions $(v_1\,{:}\,\T_1)\mapsto (v_2\,{:}\,\T_2) \mapsto \cdots \mapsto (v_n\,{:}\,\T_n) \mapsto e$ are abbreviated 
as $(v_1\,{:}\,\T_1)\cdots(v_n\,{:}\,\T_n)\mapsto e$. Argument types may be omitted when
they can be inferred from the body.

The semantics differ slightly from that of traditional functional languages: arrow types $\T_1\to\T_2$
are interpreted as {\bf mappings} from values of type $\T_1$ to values of type $\T_2$. Algebraically,
interpretations of types, $\semp{\T_1}$, $\semp{\T_2}$, are sets, and interpretations of arrow-typed terms,
$f : \T_1\to\T_2$, are {\bf partial functions} --- $\semp{f} : \semp{\T_1}\rightharpoonup\semp{\T_2}$.
This implies that a term $t : \T$ may have an \newterm{undefined} value, $\semp{t}=\bot_\T$
(We would shorten it to $\semp{t}=\bot$ when the type is either insignificant or understood from the context).
For simplicity, we shall identify $\bot_{\T_1\to\T_2}$ with the empty mapping $(v:\T_1)\mapsto\bot_{\T_2}$.

All functions are naturally extended, so that $f\,\bot=\bot$.

\subsection{Operators}
\label{lang:operators}

The core language is augmented with the following intrinsic operators:

\begin{itemize}
  \item A fixed point operator $\fix f$, with denotational semantics
    \[\semp{\fix f} ~=~ \theta \textrm{ ~ s.t. ~ } \semp f\,\theta=\theta\]
  we assume that recurrences given in specifications are well-defined, 
  such that $\semp f$ has a single fixed point.
  In other words, we ignore nonterminating computations ---
  we assume that the user provides only terminating recurrences in specifications.%\rsc{This sounds like we have a way of figuring out if a computation is non-terminating. Do we? Can we get rid of this last statement?}
  \item A guard operator $[\,]_{_\square}\,$, which comes in two flavors: 
  \[\begin{array}{l}
      [x]_{\mathit{cond}} = \begin{cases}x & \mathit{cond} \\ \bot & \lnot\mathit{cond}\end{cases} \\
      {}[f]_{P_1\times P_2\times \cdots P_n} = \overline{x} \mapsto [f\,\overline{x}]_{\bigwedge P_i(x_i)} 
    \end{array}\]
  where $\overline{x} = x_1 \cdots x_n$. This second form can be used to
  refer to quadrants of an array; in this form, it always produces a term of
  type $\square\to\mathcal{R}$, where $\mathcal{R}$ is the domain of $f$.
  \item A slash operator $/$~:     \vspace{-2pt}
  \[\renewcommand\arraystretch{1.5}
    \begin{array}{ll}
      \mbox{For scalars $x,y:\S$} & x/y = \begin{cases}x & \mbox{if }x\neq\bot \\ y & \mbox{if }x=\bot\end{cases} \\
      \mbox{For $f,g:\T_1\to\T_2$} & f/g = (v:\T_1)\mapsto (f\,v)/(g\,v)
    \end{array}\]
  This operator is typically used to combine computations done on
  parts of the array. For example, \[\psi\mapsto \big[f\,\psi\big]_{I_0} ~ \Big/ ~ \big[g\,\psi\big]_{I_1}\]
  combines a result of $f$ in the lower indices of a (one-dimensional) array
  with a result of $g$ in the higher indices ($I_0$ and $I_1$, respectively, are the index subsets).
  Notice that this does not limit the areas from which $f$ and $g$ read;
  they are free to access the entire domain of $\psi$.
\end{itemize}

In our concrete syntax, function application and $\fix$ take precedence over $/$,
and the body of $\mapsto$ spans as far as possible (like $\lambda v$ in $\lambda$-calculus).

\subsection{Primitives}

The standard library contains some common primitives:

\newcommand\forallT{\forall\T\!.~}
\newcommand\forallTS{\forall\T\S\!.~}


\begin{itemize}
  \item $\R$, a type for real numbers; $\N$ for natural numbers; $\B$ for Boolean true/false.
  \item ${=} : \forallT \T\to\T\to\B$, always interpreted as equality.
  \item ${+}, {-} : \forallT \T\to\T\to\T$, polymorphic binary operators.
  \item ${<} : \forallT \T\to\T\to\B$, a polymorphic order relation.
  \item $cons : \forallT \T\to(\N{\to}\T)\to(\N{\to}\T), nil : \forallT \N\to\T$, list constructors.
  \item $\min, \max, \Sigma : \forallTS (\T\to\S)\to\S$, reduction (aggregation) operators
    on ordered/unordered collections. The collection is represented by a mapping $f : \T\to\S$,
    so that e.g. \[\semp{\min f} = \min \big\{\semp{f}\,v \;\big|\; v\in\semp{\T}, \semp{f}\,v\neq\bot\big\}\]
    The collections are assumed to be finite.
\end{itemize}

\subsection{Additional Notation}

We also adopt some syntactic sugar to make complex terms more manageable:

\begin{itemize}
  \item $x \applt f ~=~ f\,x$ for application from the left.
  \item $\langle t_1,\cdots,t_n\rangle = cons~t_1 ~ (cons \cdots (cons~t_n ~ nil) \cdots)$
    for fixed-length lists.
\end{itemize}

\subsection{Types and Type Qualifiers}
\label{lang:types}

\cbstart%
We extend the type system with predicate abstraction in the form of logically qualified data types 
(Liquid Types~\cite{PLDI08/Rondon}). 
These are refinement types that provide a natural encoding of predicate abstraction
by restricting refinements to set of abstraction predicates,
called \newterm{qualifiers}, which are defined over the base types.
\cbend{}
Contrary to their general use, the purpose of these qualiiers in Bellmania is not to
check a program for safety and reject ill-typed programs, but rather to serve as annotations for
tactics, to convey information to the solver for use in proofs, and later to help the compiler
properly emit memory access and parallelization primitives.

More specifically, typing $f : \{v:\T_1 ~|~ P(v)\}\to\T_2$ would mean that $f\,x$
can only be defined where $P(x)$ is true; otherwise, $f\,x=\bot$. 
It {\bf does not} mean that the compiler has to prove $P(x)$ at the point where the term $f\,x$
occurs.

As such, we define a Bellmania program to be well-typed iff it is well-typed
without the annotations (in its \newterm{raw form}). Qualifiers are processed
as a separate pass to properly annotate sub-terms.

Some qualifiers are built-in, and more can defined by the user. To keep the syntax simple, we somewhat
limit the use of qualifiers, allowing only the following forms:

\begin{itemize}
  \item $\{v:\T ~|~ P(v)\}$, abbreviated as $\T\cap P$. When the signature of $P$ is known (which is
  usually the case), it is enough to write \cbstart the type as\cbend{} $P$.
  \item $\{v:\T ~|~ P(v)\land Q(v)\}$, abbreviated as $\T\cap P\cap Q$, or just $P\cap Q$. This extends
  to any number of conjuncts of the same form.
  \item $(x : \T_1) \to \{v:\T_2 ~|~ R(x,v)\} \to \T_3$, abbreviated as $\big((\T_1\times\T_2)\cap R\big)\to\T_3$.
  The qualifier argument $x$ {\bf must} be the preceding argument; this extends to predicates of
  any arity (that is, a $k$-ary predicate in a qualifier is applied to the $k$
  arguments to the left of it, including the one where it appears).
\end{itemize}


\medskip  
The type refinement constructors $\cap$ and $\times$ may be composed to create \newterm{droplets} \cbstart\diffnote{reason for the name}(tiny bits of liquid)\cbend,
using the abstract syntax in \Cref{lang:droplets}.
Note that the language does not include tuple types; hence all function types are
implicitly curried, even when using $\times$.
Droplets can express conjunctions of qualifiers,
as long as their argument sets are either disjoint or contained, but not overlapping;
for example, \[x:\{v:\T_1~|~P(v)\}\to \{v:\T_2~|~Q(v)\land R(x,v)\}\to\T_3\] can be written as
$\big((P\times Q)\cap R\big)\to\T_3$, but \[x:\T_1\to y:\{v:\T_2~|~R(x,v)\} \to \{v:\T_3~|~R(y,v)\}\to\T_4\]
cannot be represented as a droplet, except by extending the vocabulary of qualifiers.

\cbstart\diffnote{motivate droplets}%
Through the use of droplets, we manage to harness refinement types without being to 
cumbersome and requiring the user to write verbose annotations.
To this end, we trade off some expressive power for succinctness.
\cbend

\newcommand\Jsquaredlt{J^2\!\!\!{\scriptscriptstyle _{<}}}

%\exampleTitle
\begin{example}

\noindent
\cbstart\diffnote{this example has been moved up, replacing a shorter example and showcasing types inside terms}%
The specification of the Parenthesis problem (\Cref{overview:paren spec}) will be written as

\vspace{-5mm}
\[
  \renewcommand\arraystretch{1.2}
  \begin{array}{@{}l@{}l@{}l@{}}
    \lspan2{x : J \to \R} \\
    \lspan2{w : J^3 \to \R} \\
    G ~=~ \fix~(\theta:\Jsquaredlt\to\R)\,i\,j \mapsto{}
      & \big[x_i\big]_{i+1=j} ~\Big/ \\
      & \min k\mapsto\theta_{ik}+\theta_{kj}+w_{ikj}
  \end{array}
\]

$\Jsquaredlt$ is used here as an abbreviation for $(J{\times}J)\,{\cap}\,{<}$.
We also use $f_{xy}$
as a more readable typographic alternative for $f\,x\,y$,
where $f$ is a function and $x$, $y$ are its arguments.

Note that the range for $k$ in $\min k\mapsto\cdots$ is implicit, given the type of
$\theta$: 
\vspace{-.5em}
\[\theta_{ik}\neq\bot\implies i<k \quad \mbox{and} \quad \theta_{kj}\neq\bot\implies k<j\]

\end{example}

%\vspace{-.5em}
%\hrule
%\cbend
%\medskip

\iffalse
\newcommand\examplePar{%
\vspace{1pt}\noindent\hspace{-2pt}%
\tikz[baseline=(E.base)]\node(E)[rectangle,draw=black!40!white,rounded corners=.7em] {\bf Example.};
}
\examplePar
The array $G$ computed in the Parenthesis example (\Cref{overview:paren spec}) can be typed using:
\[
\begin{array}{l}
  G : ((J\times J)\cap{<}) \to \R \\
\end{array}
\]

This states that $G\,i\,j$ is only defined for $i<j$. It doesn't {\em force} it to be defined,
as it is still a partial function.
\fi

\subsubsection{Typing Rules}

\begin{figure}
\[
\begin{array}{lcll}
  d       & ::= & e^1 \quad |\quad e^k\to d \\
  e^1     & ::= & \T & \mbox{\it\small for scalar type $\T$} \\
  e^{k+l} & ::= & e^k \times e^l \\
  e^k     & ::= & e^k \cap P & \mbox{\it\small for $k$-ary predicate symbol $P$} 
\end{array}
\]
\vspace{-.5em}
\caption{\label{lang:droplets}
  Syntax of type qualifiers (droplets). $k$, $l$ are positive integers
  that stand for dimension indexes.}
\end{figure}

As mentioned earlier, annotations are ignored when typechecking a term.
This gives a simple characterization of type safety without the need to
explictly write any new typing rules. It also means that for $f:\T_1\to\T_2$, $x:\T_3$, we obtain $f\,x:\T_2$ whenever
$\T_1$ and $\T_3$ have the same \newterm{shape} \cbstart(the raw type obtained by removing all qualifiers)\cbend.
This requires some explanation.

Considering a (partial) function $\T\to\S$ to be a set of pairs of elements $\langle x,y\rangle$ 
from its domain $\T$ and range $\S$, respectively, it is clear to see that any function of type $\T_1\to\S_1$,
such that $\semp{\T_1}\subseteq\semp{\T}$, $\semp{\S_1}\subseteq\semp{\S}$, 
is \emph{also}, by definition, a function of type $\T\to\S$, since $\semp{\T_1}\times\semp{\S_1}\subseteq\semp{\T}\times\semp{\S}$.
If we define subtyping as inclusion of the domains, i.e. $\T_1 <:\T$ whenever $\semp{\T_1}\subseteq\semp{\T}$,
this translates into:
%
\[\T_1<:\T ~\land~ \S_1<:\S ~~\Rightarrow~~ (\T_1\to\S_1) <: (\T\to\S)\]

In this case, the type constructor $\to$ is {\bf covariant} in both arguments.\footnote{This is different from classical view, and holds in this case because we chose to interpret functions as \emph{mappings} (see beginning of this section).}
With this in mind, a function $g:(\T\to\S)\to \S_2$ can be called with an argument $a: \T_1\to\S_1$,
by regular subtyping rules, and $g\,a : \S_2$.

When the argument's type is not a subtype of the expected type, but has the same shape,
it is \newterm{coerced} to the required type by restricting values to the desired proper subset.
%
\[\mbox{For }h:\T\to\S \qquad \semp{h\,a} ~=~ \semp{h}\big(\semp{a} :: \T\big)\]

Where $::$ is defined as follows:
\begin{itemize}
  \item For scalar (non-arrow) type $\T$, \[x :: \T ~=~ \begin{cases}x & \mbox{if }x\in\semp{\T} \\ \bot & \mbox{if }x\not\in\semp{\T}\end{cases}\]
  \item $f :: \T\to\S ~=~ x\mapsto \big(f\,(x :: \T)\big) :: \S$
\end{itemize}

\medskip
We extend our abstract syntax with an explicit \newterm{cast operator}
$t::\T$ following this semantics. Notice that this is not the same as $t:\T$, which is a \emph{type judgement}.

\subsubsection{Type Inference}

\begin{figure*}
\vspace{-8mm}
\newcommand\typerule[2]{{\renewcommand\arraystretch{1.5}\begin{array}[t]{c} #1 \\ \hline #2 \end{array}}}
\[
\renewcommand\arraystretch{2.8}
\begin{array}{cccp{1cm}}
  \multirow{2}{1cm}[-2em]{\begin{tikzpicture}\draw (0,0) -- ++(-1mm,0) -- ++(0,-10.2em) {} node[midway,anchor=base,xshift=-3mm,rotate=90] {Core language} -- ++(1mm,0);\end{tikzpicture}} &
  \typerule{e=v \qquad \Gamma,v:\T_1\vdash e:\T_0}
           {\Gamma,v:\T_1 ~\vdash~ e:\T_0\sqcap\T_1} &
  \typerule{e=e_1\,e_2 \qquad \Gamma \vdash e:\T, ~ e_1:\T_1\to\S_1, ~ e_2:\T_2}
           {\renewcommand\arraystretch{1.2}
            \begin{array}[t]{@{}l@{}l@{}}
              \Gamma ~\vdash~ & e:\T\sqcap\S_1,  \quad
                                e_2:\T_1\sqcap\T_2, \\
                              & e_1:(\T_1\to\S_1)\sqcap(\T_2\to\T)
            \end{array}} & \\
  &
  \cspan2{
  \typerule{e=(v:\T)\mapsto e_1 \qquad \Gamma\vdash e:\T_0\to\S_0 \qquad \Gamma,v:\T\sqcap\T_0\vdash e_1:\T_1}
           {\renewcommand\arraystretch{1.2}
            \begin{array}[t]{r@{}l}
              \Gamma ~\vdash~ & e:(\T_0\to\S_0)\sqcap(\T\to\T_1)\\
              \Gamma,v:\T\sqcap\T_0 ~\vdash~ & e_1:\T_1\sqcap\S_0 
            \end{array}}  } & \\
  %
  \multirow{2}{1cm}[-1.8em]{\begin{tikzpicture}\draw (0,0) -- ++(-1mm,0) -- ++(0,-7.7em) node[midway,anchor=base,xshift=-3mm,rotate=90] {Extensions} -- ++(1mm,0);\end{tikzpicture}} &
  \typerule{e=\fix e_1 \qquad \Gamma\vdash e:\T, ~ e_1:\T_1\to\T_2}          % fix
           {\Gamma\vdash e: \T\sqcap\T_2} &
  \typerule{e=e_1/e_2 \qquad \Gamma\vdash e:\T, ~ e_1:\T_1, ~ e_2:\T_2}      % /
           {\renewcommand\arraystretch{1.2}
            \begin{array}[t]{@{}l@{}l@{}}
              \Gamma\vdash{} & e_1:\T_1\sqcap\T, ~ e_2:\T_2\sqcap\T
            \end{array}} & \\
  &
  \typerule{e=[e_1]_{_{cond}} \qquad \Gamma\vdash e:\T,e_1:\T_1}                           % ::
           {\Gamma\vdash e:\T\sqcap\T_1, ~ e_1:\T\sqcap\T_1} &
  \typerule{e=e_1::\T \qquad \Gamma\vdash e:\T_0,e_1:\T_1}                           % ::
           {\Gamma\vdash e:\T\sqcap\T_0\sqcap\T_1, ~ e_1:\T\sqcap\T_0\sqcap\T_1}
\end{array}
\]
\caption{\label{lang:type refinement rules}
  Type refinement rules, for inferring qualifiers in sub-expressions.}
  \vspace{2mm}
\end{figure*}

Base types are inferred normally as in a classical Hindley-Milner type system~\cite{78/Milner}.
The operators (\Cref{lang:operators}) behave like polymorphic
constants with the following types:
\[\renewcommand\arraystretch{1.5}
  \begin{array}{c}
    {\fix} ~:~ \forallT ~(\T\to\T)\to\T \qquad 
    {/} ~:~ \forallT ~\T\to\T\to\T \\
    (::\T) ~:~ \mathrm{shape}[\T]\to\mathrm{shape}[\T]
  \end{array}\]

As for $[\,]_{_\square}$, for all typing purposes the first variant is a no-op, and the second variant
is just syntactic sugar for $:: \square\to\_$\,,
where \,$\_$\, is a fresh type variable.

Any type variables occurring in type expressions are resolved at that
stage through unification. In particular, it means that type variables are always assigned raw types.

Qualifiers are also inferred by propagating them up and down the syntax tree.
Since the program already typechecks once the base types are in place, the problem is no longer
one of finding {\em valid} annotations, but rather of {\em tightening} them as much as possible
without introducing semantics-changing coercions. For example, the term $(f :: I\to(I\cap P))\,i$ may
be assigned the type $I$, but it can also be assigned $I\cap P$ without changing its semantics.

Qualifiers are propagated by defining a \newterm{type intersection} operator $\sqcap$ that
takes two droplets of the same shape $\T_1$, $\T_2$ and returns a droplet with a conjunction of all the qualifiers
occuring in either $\T_1$ or $\T_2$. The operator is defined in terms of the corresponding liquid types:
\begin{itemize}
  \item If $\T_1=\{v:\T ~|~ \varphi_1\}$ and $\T_2=\{v:\T ~|~ \varphi_2\}$,
	\[\T_1\sqcap\T_2 ~=~ \{v:\T ~|~ \varphi_1\land\varphi_2\}\]
  \item If $\T_1=x{\,:\,}\S_1\to\S_2$, $\T_2=x{\,:\,}\S_3\to\S_4$ (normalized so that $\T_1$ and $\T_2$ use the same names for arguments),
    \[\T_1\sqcap\T_2 ~=~ x{\,:\,}(\S_1\sqcap\S_3)\to(\S_2\sqcap\S_4)\]
\end{itemize}

We then define the \newterm{type refinement} steps for terms. They are listed in \Cref{lang:type refinement rules}.
These rules are applied continuously until a fixed point is reached.
The resulting types are eventually converted back to droplet form (expressed via $\cap$ and $\times$);
qualifiers that cannot be expressed in droplets are discarded.

\medskip
\cbstart\diffnote{added some justification for the type refinement rules.}%
The top three rules apply to ordinary constructs of typed $\lambda$-calculus:
the first rule propagates qualifiers from the environment to leaf expressions that use a declared variable;
the next two rules propagate information across application and refinement terms.
For an application $f\,i$, both the type of $i$ and the domain of $f$ can be refined
via the other, e.g. if $f:I_0\to\R$ and $i:I$, then $i$ can be refined to $I_0$; 
symmetrically, if $f:I\to\R$ and $i:I_0$, then $f$ can be refined to $I_0\to\R$.
Similarly, the range of $f$ and the type of $f\,i$ can refine each other.
The rule for abstraction terms is analogous, except that information about the type
of the argument $v$ is passed through the environment.

The bottom rules pertain to language extensions defined previously in this section.
For $/$, qualifiers can safely seep down to sub-expressions (but not back up), e.g. if $x/y:I_0$,
it is safe to assert $x:I_0$, $y:I_0$ without changing the semantics of the expression.
For casts, $x$ and $x::\T$ are tied to have the same refinements, and those
of $\T$ are appended. 
A condition $[x]_{_\square}\,$ ties the sub-term type in the same way, but without
adding any qualifiers. The rule for $\fix$ is the only one that requires some thought:
the type of $\fix f$ cannot be forced down on $f$; but since, by definition, $\fix f = f\,(\fix f)$
must hold, then any restrictions imposed by the range of $f$ can be asserted for $\fix f$.
\cbend

Note that two syntactically identical terms in different sub-trees may be assigned
different types by this method. This is a desirable property, as (some) context information
gets encoded in the type that way.

The interested reader can find an example of how type inference is applied in \Cref{annex:example:type inference}.

\begin{comment}
\medskip
\examplePar Say $I$, $T$ are types, $\widehat{I}_0 : I\to\B$ a qualifier, $\S$ a type variable, $0:T$ a constant.
The term $(f:I_0\to\S)\,i\mapsto f\,i\,i ~/~ 0$ will induce the following type
inferences:

\begin{tikzpicture}[node distance=1em]
  \node(farg) {$(f:I_0\to\S)$};
  \node(iarg)[right=of farg] {$i$};
  \node(mapsto)[right=of iarg] {$\mapsto$};
  \node(f)[right=of mapsto] {$f$};
  \node(i1)[right=of f] {$i$};
  \node(i2)[right=of i1] {$i$};
  \node(slash)[right=of i2] {$\big/$};
  \node(zero)[right=of slash] {$0$};
  \node(l0)[coordinate,below=1mm of farg] {};
  \node(l1)[coordinate,below=4.5mm of l0] {};
  \node(l2)[coordinate,below=3.5mm of l1] {};
  \node(l3)[coordinate,below=3.5mm of l2] {};
  \node(l4)[coordinate,below=3.5mm of l3] {};
  \def\ytip{2pt}
  \draw (l0 -| farg.west) -- +(0,-\ytip) -| (l0 -| farg.east) node[pos=0.25,below] {\tiny $I_0\to I\to T$};
  \draw (l0 -| iarg.west) -- +(0,-\ytip) -| (l0 -| iarg.east) node[pos=0.25,below] {\tiny $I$};
  \draw (l0 -| f.west) -- +(0,-\ytip) -| (l0 -| f.east) node[pos=0.5,anchor=23,inner xsep=0] {\tiny $I_0\to I\to T$};
  \draw (l0 -| i1.west) -- +(0,-\ytip) -| (l0 -| i1.east) node[pos=0.25,below] {\tiny $I_0$};
  \draw (l0 -| i2.west) -- +(0,-\ytip) -| (l0 -| i2.east) node[pos=0.25,below] {\tiny $I$};
  \draw (l0 -| zero.west) -- +(0,-\ytip) -| (l0 -| zero.east) node[pos=0.25,below] {\tiny $T$};
  \draw (l1 -| f.west) -- +(0,-\ytip) -| (l1 -| i1.east) node[pos=0.25,below] {\tiny $I_0\to T$};
  \draw (l2 -| f.west) -- +(0,-\ytip) -| (l2 -| i2.east) node[pos=0.25,below] {\tiny $T$};
  \draw (l3 -| f.west) -- +(0,-\ytip) -| (l3 -| zero.east) node[pos=0.25,below] {\tiny $T$};
  \draw (l4 -| farg.west) -- +(0,-\ytip) -| (l4 -| zero.east) node[pos=0.25,below] {\tiny $(I_0\to I\to T)\to I\to T$};
\end{tikzpicture}
%\hspace*{\fill}

\noindent
In this case, the type variable $\S$ has been assigned $I\to T$.
\end{comment}


\section{Tactics}
\label{tactics}

We now define the method with which our framework transforms program terms, by means of \newterm{tactics}.
A tactic is a scheme of equalities that can be used for rewriting.
When applied to a program term, any occurrence of the {\bf left-hand side} is replaced by the {\bf right-hand side}.%
\footnote{This is also a standard convention in Coq~\cite{COQ}, for example.}
A valid application of a tactic is an instance of the scheme that is well-typed and logically valid
(that is, the two sides have the same interpretation in any structure that interprets the free
variables occurring in the equality).

The application of tactics yields a sequence of program terms, each of which is checked to
be equivalent to the previous one. We refer to this sequence by the name \newterm{development}.

We associate with each tactic some \newterm{proof obligations}, listed after the word \textbf{\textit{Obligations}}
in the following paragraphs.
When applying a tactic instance, these obligations are also instantiated and given to an automated prover. 
If verified successfully, they entail the validity of the instance.
\cbstart\diffnoteup{\vspace{-.5cm}a note about incompleteness, and also where the obligations come from}%
In some cases, the obligations are strictly stronger than the validity of the rewrite,
which means the TAE is not complete and may fail to apply a valid transformation.
It is important to note that the user does not have to specify, or in fact be aware of,
any proof obligations; they are defined as part of the tactic library.
\cbend

The following are the major tactics provided by our framework. 
More tactic definitions are given in \Cref{annex:more tactics}.

\newcommand\Obligations{\medskip\noindent\textbf{\textit{Obligations}:} }
\newcommand\reduce{\operatorname{reduce}}
\newcommand\listConcat{{\scriptstyle \,++\,}}

\theoremstyle{definition}
\newtheorem{tactic}{Tactic}

\newcommand\tacticdef[1]{\subsection*{\sf\larger #1}}
\newcommand\tacticdefcompact[1]{\medskip\noindent{\sf\larger #1}\medskip\hfill}

\tacticdefcompact{Slice} \label{tactics:Slice}
$f ~=~ \big[f\big]_{X_1} ~\Big/~ \big[f\big]_{X_2} ~\Big/ ~\cdots~ \Big/~ \big[f\big]_{X_r}$\hspace{2mm}

This tactic partitions a mapping into sub-regions. Each $X_i$ may be a cross product ($\times$)
according to the arity of $f$.%
\cbstart\diffnote{example for $X_i$'s}
For example, $X_1=J_0\times J_0$, $X_2=J_0\times J_1$, $X_3=J_1\times J_1$.
\cbend

\Obligations \cbstart\diffnote{removed the confusing "Obligations: tactic."}just the equality above.
\cbend

Informally, the recombination expression is equal to $f$
when $X_{1..r}$ ``cover'' all the defined points of $f$ (also known as the \newterm{support} of $f$).

\medskip
\tacticdefcompact{Stratify}
$\fix (f\applt g) ~=~ (\fix f) ~\applt~ \big(\psi\mapsto \fix (\widehat\psi\applt g)\big)$
\\
where $\widehat\psi$ abbreviates $\theta\mapsto\psi$, with fresh variable $\theta$.

\cbstart\diffnote{made this clearer with a concrete example and an illustration}%
This tactic is used to break a long (recursive) computation into steps. $\psi$ is bound the result of the first computation ($f$), and then used as the input for the second step ($g$).
Typically, we apply \textsf{Stratify} to a computation that has previously been \textsf{Slice}d (although that does not have to be the case); to understand the efficacy of \textsf{Stratify}, we present a minimal example.

Suppose the program $\fix \theta{\,\mapsto\,}\big([t]_{I_0}\big/[t]_{I_1}\big)$, where $t=i\mapsto\theta_{i-1}+1$.
We would like to obtain a program that computes the lower half $I_0$ first, then the upper half $I_1$; that is, 
$\big(\fix\theta\mapsto[t]_{I_0}\big/[\psi]_{I_1}\big) \applt
  \big(\psi \mapsto \fix\theta\mapsto[\psi]_{I_0}\big/[t]_{I_1}\big)$.

\begin{center}
\begin{tikzpicture}[node distance=.5cm, >=latex,
  array/.style={rectangle,draw,
                minimum width=2cm, minimum height=5mm},
  div/.style={shorten <=.5pt,shorten >=.5pt},
  block/.style={rectangle,fill=Orange, fill opacity=0.1, inner sep=0pt},
  src/.style={rounded corners=3pt,draw,dotted,fill=Orange, fill opacity=0.1, inner sep=1.5pt}]
  \node(s0)[array] {};
  \draw[div] (s0.north) -- (s0.south);
  
  \node(applt)[right=of s0] {$\applt$};
  
  \node(s1)[array,right=of applt] {};
  \draw[div] (s1.north) -- (s1.south);
  
  % labels
  \node(s0lbl)[below=0mm of s0] {$\fix\theta\mapsto[t]_{I_0}\big/[\psi]_{I_1}$};
  \node(s1lbl)[below=0mm of s1,anchor=140] {$\psi \mapsto \fix\theta\mapsto[\psi]_{I_0}\big/[t]_{I_1}$};
  
  % decoration
  \node(s0blk)[block,src,fit={(s0.north west) (s0.south)}] {};
  \node(s1blk)[block,fit={(s1.north) (s1.south east)}] {};
  
  \draw (s0blk.20) edge[->,out=10,in=160] (s1blk.center);
\end{tikzpicture}
\end{center}

This can be obtained by choosing 
$f=\theta\mapsto[t]_{I_0}\big/[\psi]_{I_1}$
and
$g=f\,\theta \mapsto [f\,\theta]_{I_0}\big/[t]_{I_1}$.

Notice how $f\,\theta$ is used as a placeholder for the sub-computation $f$ in $g$.
On the one hand, $f\applt g = \theta\mapsto\big[[t]_{I_0}\big/[\psi]_{I_1}\big]_{I_0}\big/[t]_{I_1}$, which,
modulo some simplification, is equivalent to the original term $\theta\mapsto[t]_{I_0}\big/[t]_{I_1}$.
On the other hand, when instantiating \textsf{Stratify}, $f\,\theta$ becomes $\widehat\psi\,\theta$, 
which is equivalent to $\psi$, giving the term for the second phase.
\cbend

$\psi$ may be fresh, or it may reuse a variable already occurring in $g$, rebinding those occurrences.
This precisely encodes our intuition of computing the first fixed point \textit{in situ},
then the second one based on the result of the first.

\Obligations Let $h=f\applt g$ and $g'=\psi\mapsto\widehat\psi\applt g$. Let $\theta,\zeta$ be
fresh variables; then,\vspace{-.3em}
\begin{equation}
\renewcommand\arraystretch{1.5}
\begin{array}{l@{\qquad}l}
f\,(g'\,\zeta\,\theta) ~=~ f\,\zeta &
g'\,(f\,\theta)\,\theta ~=~ h\,\theta
\end{array}
\label{tactics:Stratify obligations}
\end{equation}

%Proof is given in \Cref{tactics:soundness}.

\vspace{.1em}
\tacticdefcompact{Synth} \label{tactics:Synth}
$\fix\big(h_1 ~\big/~ \cdots ~\big/~ h_r\big) ~=~ 
  f_1 :: \T_1 ~\big/~ \cdots ~\big/~ f_r :: \T_r$

This tactic is used to generate recursive calls to sub-programs. For $i=1..r$, $f_i$
is one of the following: $\fix h_i$, $h_i\,\psi$, or $t\,\psi$, where $\psi$ is some
variable and $t$ is a term corresponding to a previously defined subroutine
($A$, $B$, $C$ in the example).
Bellmania chooses these values automatically (see \Cref{tactics:synthesis}),
but the user may override it.

\newcommand\Y{\mathcal{Y}}

\Obligations Let $h=h_1/\cdots/h_r$, and let $\T\to\T$ be the shape of $h$. 
  For each $f_i$, depending on the form of $f_i$:
\begin{itemize}
  \item If $f_i \cong \fix f$ (for some $f$) --- \\
    \rule{0pt}{12pt}
    $h :: (\T \to \Y) = h :: (\Y \to \Y) = f :: (\Y \to \T)$
    for some $\Y$ which is a subtype of $\T$ and a supertype of $\T_i$.
  \item If $f_i$ does not contain any ``$\fix$'' terms ---\\
    \rule{0pt}{12pt}
    $h\,(h\,\theta) :: \T_i = f_i :: \T_i$ for a fresh variable $\theta$.
\end{itemize}

$\cong$ denotes syntactic congruence up to $\beta$-reduction.

\bigskip
\begin{theorem}
Let $s=s'$ be an instance of one of the tactics introduced in this section.
let $a_i=b_i$, $i=1..k$, be the proof obligations. If $\semp{a_i}=\semp{b_i}$
for all interpretations of the free variables of $a_i$ and $b_i$, then
$\semp{s}=\semp{s'}$ for all interpretations of the free variables of $s$ and $s'$.
\end{theorem}

Proof is given in \Cref{annex:proof}.


\newenvironment{tacticbox}[1]{\begin{center}
  \begin{tabular}{|@{~~~~}l@{~~~~}|}\hline
    \rule{0pt}{2.3ex}\underline{\sf \,#1\,}\\[.4em]$}
  {$\\[-1em] \\[.3ex] \hline \end{tabular} \end{center}}

\newenvironment{tacticboxInline}[1]{\begin{tabular}{|@{~~~}l@{~~~}|}\hline
    \rule{0pt}{2.3ex}\underline{\sf \,#1\,}\\[.4em]$}
  {$\\[-1em] \\[.3ex] \hline \end{tabular}}

\newcommand\vtyped[2]{\underset{\scriptscriptstyle ( #2 )}{ #1 }}
\newcommand\htyped[2]{({#1} : {\scriptstyle{#2}})}
\newcommand\htypednp[2]{{#1} : {\scriptstyle{#2}}}

% \quadrants{A}{B}{C}{D}   -- normal sized quadrants
% \quadrants*{A}{B}{C}{D}  -- small quadrants
\makeatletter
\newcommand{\quadrants@normal}[4]{
  \renewcommand\arraystretch{1.5}
   \begin{array}{c|c}
     #1 & #2 \\ \hline
     #3 & #4
   \end{array}}
\newcommand{\quadrants@small}[4]{
  \renewcommand\arraystretch{0.9}
   \begin{array}{@{~}c@{~}|@{~}c@{~}}
     \scriptstyle #1 & \scriptstyle #2 \\ \hline
     \scriptstyle #3 & \scriptstyle #4
   \end{array}}
\newcommand\quadrants{\@ifstar\quadrants@small\quadrants@normal}
\makeatother


%\exampleTitle
\begin{example}

\noindent
The na\"ive implementation of \Cref{intro:naive}
can be written as
%
\newcommand\Jsquaredltsml{J^2_{\mbox{\fontsize{6}{6}\selectfont \hspace{.5pt}<}}}
\begin{equation}
  \renewcommand\arraystretch{1.2}
  \begin{array}{@{}r@{}l@{}}
    \Psi ~=~ & i\,j\mapsto
	           [x_i]_{i+1=j} \\[2mm]
%
    A^{^{J}} ~=~ 
	      & \psi\mapsto \fix 
	        \htyped{\theta}{\Jsquaredltsml\to\R}~ \htyped{i}{J} ~ \htyped{j}{J}
	        ~\mapsto \\
	      &
          ~~
	      \begin{array}{@{}l@{}l@{}} 
	        \min\,\langle~ & \psi_{ij}, \\
	         & \min\, \htyped k J \mapsto\theta_{ik}+\theta_{kj}+w_{ikj}
	        ~\rangle
	      \end{array}
  \end{array}
  \label{tactics:paren phase A}
\end{equation}

%Vertical typeset was used to save some horizontal space, but $\vtyped v\T$
%should be read as just $v:\T$.

As mentioned in \Cref{overview}, the first step Richard does is to apply
{\sf Slice} to the program, thus introducing a partitioning into quadrants.

\smallskip

\begin{tacticbox}{Slice}
  \begin{array}{@{} l  l @{}}
       f ~=~ \theta\,i\,j\mapsto \cdots \\
       X_1 ~=~ \_\times J_0\times J_0 &
       X_2 ~=~ \_\times J_0\times J_1 \\
	   &
       X_3 ~=~ \_\times J_1\times J_1 \\[.5em]
       \cspan2{\mbox{\small ({\it each} ``\_'' {\it is a fresh type variable})}}
  \end{array}
\end{tacticbox}
\vspace{-2mm}
%
\begin{flalign}
  ~A^{^J} = \psi\mapsto \fix \quadrants{\qbox1}{\qbox2}{}{\qbox4} &&
  \label{tactics:A sliced}
\end{flalign}
\vspace{-2mm}
\[
  \renewcommand\arraystretch{1.2}
  ~~~
  \begin{array}{@{}r@{}l@{~}c@{~}c@{}l@{}l@{}}
	\qbox1 ~=~ & \theta\,& \htyped{i}{J_0} & \htyped{j}{J_0} & {}\mapsto  \\ 
	 & \lspan4{\qquad \min\,\langle~ \psi_{ij}, \min\, \htyped k J \mapsto\theta_{ik} + \theta_{kj} + w_{ikj} ~\rangle} \\
	\qbox2 ~=~ & \theta\,& \htyped{i}{J_0} & \htyped{j}{J_1} & {}\mapsto \\
	 & \lspan4{\qquad \min\,\langle~ \psi_{ij}, \min\, \htyped k J \mapsto\theta_{ik} + \theta_{kj} + w_{ikj} ~\rangle} \\
	\qbox4 ~=~ & \theta\,& \htyped{i}{J_1} & \htyped{j}{J_1} & {}\mapsto \\
	 & \lspan4{\qquad \min\,\langle~ \psi_{ij}, \min\, \htyped k J \mapsto\theta_{ik} + \theta_{kj} + w_{ikj} ~\rangle} \\
  \end{array}
\]

With repeated applications of {\sf Slice}, a program may grow to become
quite large; to make large program terms easy to read and refer to, we provide
boxed numerals \cbstart\qbox1, \qbox2, etc.\cbend{} as labels for sub-terms, using them as abbreviations for these terms where they
occur in the containing expression.

In addition, to allude to the reader's intuition, expressions of the form
$a/b/c/d$ will be written as $\quadrants*{a}{b}{c}{d}$ when the slices
represent quadrants. \cbstart The types of the quadrants should be clear from the context;
here, their types are\cbend\diffnoteup{\vspace{-1cm}emphasized the role of boxed literals in the development}%
\vspace{-.5em}
\[\htypednp{\qbox1}{\_{\times}J_0{\times}J_0\to\R},~ \htypednp{\qbox2}{\_{\times}J_0{\times}J_1\to\R},~ \htypednp{\qbox4}{\_{\times}J_1{\times}J_1\to\R}\]%
%

The type of $\theta$ in \Cref{tactics:A sliced} is still
$\Jsquaredlt\to\R$. In order to avoid too much clutter caused by type terms,
the Bellmania UI uses a heuristic and only displays some of them. By hovering,
the user can inspect types that are not shown.

The result is a functional representation of \Cref{overview:logical-slice-stratify}$(a)$;
with the added term $\psi_{ij}$ it allows $A$ to accept an input array as an argument
and minimize values that are already stored in it; not quite crucial for this particular
instance where the input is just $\Psi$, but very useful in cases where it is necessary
to split a minimization operation into several sub-ranges, to achieve the desired
memory locality.

\newcommand\lifted[1]{\widehat{#1}}

\begin{tacticbox}{Stratify \qbox 1}
  \begin{array}{@{} l @{} l @{}}
    f ~=~ \quadrants*{\tinyqbox1}{\lifted\psi}{\lifted\psi}{\lifted\psi}
         & \mbox{\small ({\it recall that } $\widehat\psi=\theta\mapsto\psi$)} \\
    g ~=~ z\mapsto\quadrants*{z}{\tinyqbox{2}}{}{\tinyqbox4} &
    \qquad\quad\psi=\psi
  \end{array}
\end{tacticbox}

\begin{equation}
  \renewcommand\arraystretch{1.5}
  \begin{array}{@{}l@{}} %@{}l@{}c@{}c@{}l@{}l@{}}
    A^{^J} =~ \psi\mapsto \left(\!\fix\, \quadrants{\qbox1}{\lifted\psi}{\lifted\psi}{\lifted\psi}\right) \applt 
       ~~\psi \mapsto \fix\, \quadrants{\lifted\psi}{\qbox2}{}{\qbox4} \\
    \quad
    \qbox1 \,\qbox2 \,\qbox4 \textit{ as in \Cref{tactics:A sliced}}
  \end{array}
  \label{evaluation:A stratified}
\end{equation}

The reason for this particular choice of $f$ and $g$ is as explained in
\Cref{tactics}. Richard does not have to worry too much about them, because
they are encoded in the tactic application engine,
so that Bellmania knows how to build them automatically
based on the term being stratified (\qbox1 in this case).

Notice that an existing variable $\psi$ is reused, rebinding any occurrences within $\qbox2$\,, $\qbox4$\,.
This effect is desirable, as it limits the context of the expression: the inner $\psi$ shadows the outer $\psi$,
meaning $\qbox2$\,, $\qbox4$ do not need to access the data that was input to $\qbox1$\,, only its
output; therefore $\qbox1$ can be computed in-place.
The proof obligations for {\sf Stratify} make sure this transition is valid.

At this point Richard can either do another {\sf Stratify} or a {\sf Synth}.
The order is insignificant, but to be consistent with \Cref{overview:slice-stratify-synth},
let us assume he chooses the former.

\medskip

\begin{tacticbox}{Stratify \qbox1}
  \begin{array}{@{} l @{} l @{}}
    f ~=~ \quadrants*{\tinyqbox1}{\lifted\psi}{\lifted\psi}{\lifted\psi}
         & \mbox{\small ({\it recall that } $\widehat\psi=\theta\mapsto\psi$)} \\
    g ~=~ z\mapsto\quadrants*{z}{\tinyqbox2}{}{\tinyqbox4} &
    \qquad\quad\psi=\psi
  \end{array}
\end{tacticbox}

\begin{equation}
  \renewcommand\arraystretch{1.5}
  \begin{array}{@{} r @{} l @{}}
    A^{^J}\! = \psi\mapsto{} & \left(\!\fix\, \quadrants{\qbox1}{\lifted\psi}{\lifted\psi}{\lifted\psi}\right) \applt 
       ~\psi \mapsto \left(\!\fix\, \quadrants{\lifted\psi}{\qbox2}{\lifted\psi}{\lifted\psi}\right) \applt \\
       & \psi \mapsto \fix\, \quadrants{\lifted\psi}{\lifted\psi}{}{\qbox4} \\
    \lspan2{
     \quad
     \qbox1 \,\qbox2 \,\qbox4 \textit{ as in \Cref{tactics:A sliced}}
    }
  \end{array}
  \hspace{-8mm}
  \label{tactics:A stratified twice}
\end{equation}

\begin{center}
\begin{tacticboxInline}{Synth \qbox1}
    \renewcommand\arraystretch{1.3}
	\begin{array}{@{} l @{\quad} l @{}}
       h_1= \qbox1            &  h_{2,3,4}=\lifted\psi  \\
	   f_1 = A^{^{J_0}}\,\psi &  f_{2,3,4} = \psi       \\
	   \cspan2{\Y = J_0^2 \to \R}
   \end{array}
\end{tacticboxInline}
%
\begin{tacticboxInline}{Synth \qbox4}
    \renewcommand\arraystretch{1.3}
	\begin{array}{@{} l @{\quad} l @{}}
       h_{1,2} = \lifted\psi   &  h_3 = \qbox4         \\
	   f_{1,2} = \psi          &  f_3 = A^{^{J_1}}\,\psi  \\
	   \cspan2{\Y = J_1^2 \to \R}
   \end{array}
\end{tacticboxInline}
\end{center}

\begin{equation}
  \renewcommand\arraystretch{1.5}
  \begin{array}{@{} r @{} l @{}}
    A^{^J}\! = \psi\mapsto{} & \quadrants{A^{^{J_0}}\,\psi}{\psi}{\psi}{\psi} ~ \applt 
       ~\psi \mapsto \left(\!\fix\, \quadrants{\lifted\psi}{\qbox2}{\lifted\psi}{\lifted\psi}\right) \applt \\[5mm]
       & \psi \mapsto \quadrants{\psi}{\psi}{}{A^{^{J_1}}\,\psi} \\
    \lspan2{
     \quad
     \qbox2 \textit{ as in \Cref{tactics:A sliced}}
    }
  \end{array}
  \hspace{-5mm}
  \label{tactics:A synthed}
\end{equation}

For $\qbox2$\,, the situation is slightly more complicated because no instance of the routine
$A$ matches the specification and there are no other routines to choose from.
Richard defines a new routine $B^{J_0J_1}$ by carving the respective
subexpression from the developed program $A^J$.
Notice that $B$ has two parameters, because it depends not only on the index
range, but also on the particular partitioning.
Next, Richard will carry on developing $B$ in a similar manner.

\end{example}
%\medskip
%\hrule
%\medskip

\subsection{Synthesis-powered {\sf Synth} Tactic}
\label{tactics:synthesis}

As mentioned in \Cref{intro,overview}, the user is assisted by automatic
inference while applying tactics. In particular, the {\sf Synth} tactic requires
the user to specify a subroutine to call and parameters to call it with.
In addition, the subtype $\Y$ is required to complete the correctness proof.
To automate this task, Bellmania employs Counterexample-guided Inductive Synthesis ({\cegis}), a software synthesis technique
implemented in the tool {\Sketch}~\cite{STTT13/Solar-Lezama}. The proof obligations, along with the possible
space of parameter assignments taken from the set of sub-types defined during
{\sf Slice}, are translated to {\Sketch}. Since {\Sketch} uses bounded domains,
the result is then verified using full SMT.

In addition to considering explicitly defined sub-types, the synthesizer also tries
small variations of them to cover corner cases. When index arithmetic is used,
the range for a sub-call may have to be extended by a row or column on one or more sides.
For each index sub-type $\T\subseteq J$, Bellmania also tries $\T\cup(\T\pm 1)$ for filling in values of parameters:
\[\T+1=\{i+1\;|\;i\in \T\} \qquad \T-1=\{i-1\;|\;i\in \T\}\]

While the number of combinations is not huge, 
it is usually hard for the user to figure out which exact call
should be made. Since {\sf Synth} is used extensively throughout the development,
This kind of automation greatly improves overall usability.

{\Sketch} times for the running example range 15--45 seconds.
For comparison, covering the same search space for a typical invocation via exhaustive search in C++ took about 1$\frac{1}{2}$ hours.

\begin{comment}
% This part I moved to the appendix

\subsection{Soundness}
\label{tactics:soundness}

\renewenvironment{proof}{\noindent{\bf Proof.~}}{}

\begin{theorem}
Let $s=s'$ be an instance of one of the tactics introduced in this section.
let $a_i=b_i$, $i=1..k$, be the proof obligations. If $\semp{a_i}=\semp{b_i}$
for all interpretations of the free variables of $a_i$ and $b_i$, then
$\semp{s}=\semp{s'}$ for all interpretations of the free variables of $s$ and $s'$.
\end{theorem}

\begin{proof}
For the tactics with \textbf{\textit{Obligations}:} tactic, the theorem is trivial.

\medskip
\noindent
{\tt >} For {\sf Stratify}, let $f$, $g$ be partial functions such that
\vspace{-.5em}
\[\renewcommand\arraystretch{1.3}
  \forall \theta,\zeta.\quad \begin{array}{l}f\,(g\,\zeta\,\theta) ~=~ f\,\zeta \quad\land\quad
  g\,(f\,\theta)\,\theta ~=~ h\,\theta
  \end{array}\quad\]
  
Assume that $\zeta = \fix f$ and $\theta = \fix (g\,\zeta)$. 
That is, $f\,\zeta = \zeta$ and $g\,\zeta\,\theta = \theta$.
Then ---
\vspace{-.5em}
\[\renewcommand\arraystretch{1.3}
  \begin{array}{l@{}l}
   h\,\theta & {}= g\,(f\,\theta)\,\theta = g\,(f\,(g\,\zeta\,\theta))\,\theta =
              g\,(f\,\zeta)\,\theta = \theta
  \end{array}\]
  
\noindent
So $\theta = \fix h$. We get $\fix h = \fix \big(g \,(\fix f)\big)$; equivalently,
\[\fix h = (\fix f) \applt \big(\psi\mapsto\fix (g\,\psi)\big)\]

Now instantiate $h$, $f$, and $g$, with $f\applt g$, $f$, and $g'$ from \Cref{tactics:Stratify obligations},
and we obtain the equality in the tactic.

\medskip
\noindent
{\tt >} For {\sf Synth}, ({\it i}) assume $f_i=\fix g$ and
\[h :: \T\to\Y = h :: \Y\to\Y = g :: \Y\to\T\]

Intuitively, $\Y$ ``cuts out'' a region of an array $\theta :: \T$ given
as input to $h$ and $g$. This area is self-contained, in the sense that
only elements in $\Y$ are needed to compute elements in $\Y$, as indicated
by the refined type $\Y\to\Y$.

Notice that from the premise follows $g :: \Y\to\T = g :: \Y\to\Y$. We use the following corollary:

\medskip\noindent
{\bf Corollary.~} Let $f : \T\to\T$; if either $f :: \T\to\Y = f :: \Y\to\Y$ or $f :: \Y\to\T = f :: \Y\to\Y$, 
then $(\fix f)::\Y = \fix (f :: \Y\to\Y)$.

Proof is included in the appendix.

\medskip
From the corollary, and for the given $h$ and $g$, we learn that $(\fix h)::\Y = \fix(h::\Y\to\Y)$,
and also $(\fix g)::\Y=\fix(g::\Y\to\Y)$. Since $h::\Y\to\Y=g::\Y\to\Y$,
we get $(\fix h)::\Y = (\fix g)::\Y$; now, $\Y$ is a supertype of $\T_i$, so $(\theta::\Y)::\T_i=\theta::\T_i$:
\[\renewcommand\arraystretch{1.3}
  \begin{array}{l}(\fix h)::\T_i = ((\fix h)::\Y)::\T_i = ((\fix g)::\Y)::\T_i = \\
    \qquad = (\fix g)::\T_i=f_i::\T_i
  \end{array}\]

{\it (ii)} Assume $h\,(h\,\theta) :: \T_i = f_i :: \T_i$ holds for any $\theta:\T$,
then in particlar, for $\theta=\fix\,h$, we get $h\,(h\,\fix h) :: \T_i = f_i :: \T_i$.
Since $h\,(h\,\fix h) = \fix h$, we obtain the conjecture $(\fix h) :: \T_i = f_i :: \T_i$.
\qed
\end{proof}

\medskip
Our reliance on the termination of $\fix$ expressions may seem conspicuous, since some of these
expressions are generated automatically by the system. However, a closer look reveals that whenever
such a computation is introduced, the set of the recursive calls it makes is a subset of those made by the existing one.
Therefore, if the original recurrence terminates, so does the new one. In any case, all the recurrences
in our development have a trivial termination argument (the indexes $i$,$j$ change monotonically between calls),
so practically, this should never become a problem.

\end{comment}
\input{05-automated}
\section{Code Generation}
\label{codegen}

We built a compiler for programs in Bellmania language that generates efficient C++ code parallelized with Intel Cilk constructs. The compiler uses type information from the development to improve the
quality of generated code. From the types of sub-terms corresponding
to array indices, the compiler extracts information about what region
of the array is read by each computation, and from the types of $\lambda$-bound index variable it constructs loops that write to
the appropriate regions. The compiler utilizes the SMT solver to infer dependency constraints as in~\cite{JACM67/Karp},
and figures out the direction of each loop (ascending or descending). 
In addition, inter-quadrant dependencies can be used to determine which computations can be run in parallel (at the granularity of
function calls) based on a fork-join model;
two calls are considered non-conflicting if each write region is disjoint from the others' read and write regions.
\cbstart\diffnote{explained the role of types in our static analysis}%
Type information is heavily utilized here: occurrences of array variables signify reads,
so their types denote (an over-approximation of) the region being ``read from'',
whereas return types denote the region being ``written to''.
Disjointness can be decided using propositional logic through the predicate abstraction induced by type qualifiers.
\cbend

The compiler also employs more traditional optimization techniques
to further simplify the code and improve its running time:
(1) lifts conditionals to loop bounds via simple interval analysis, %to optimize the loop bounds and performs some traditional compiler %transformations to reduce the number of iterations and comparisons,
(2) eliminates redundant iterations and comparisons that can be resolved at compile time,
(3) identifies loops that read non-contiguous memory blocks and applies copy optimization~\cite{ASPLOS91/Lam} automatically to better utilize caches. Examples of generated code are included in the Bellmania repository~\cite{bellmaniarepo},
and
\cbstart
(4) follows a simple heuristic that inserts loop vectorization directives.
\cbend
 
%Use Cilk for parallelization - references later
%Using SMT solver and type information to find directions of loops for fixed point computations
%Using type information and SMT solvers to parallelize different components
%Using type information to optimize the loop bounds, reduce checking of guards and do copy optimization
%Traditional compiler transformations

%code can be found in repo
%JACM67: Recurrences to loops: http://dl.acm.org/citation.cfm?id=321418
%ASPLOS1991: copy optimization: http://suif.stanford.edu/papers/lam-asplos91.pdf
\section{Empirical Evaluation}
\label{evaluation}

We implemented our technique and used it to generate ca\-che-oblivious
divide-and-conquer implementations of three algorithms that were used as
benchmarks in \cite{IPDPS15/Tithi}, and a few others.

\begin{figure}[t!]
\vspace{-5mm}
\[
  \renewcommand\arraystretch{1.2}
  \begin{array}{@{}l@{}l@{}l@{}}
    \lspan3{w^x : ((I\times I)\cap{<})\to J\to\R} \\
    \lspan3{w^y : ((J\times J)\cap{<})\to I\to\R} \\
    G ~=~ \fix \theta\,i\,j\mapsto{}
      & \lspan2{\big[0\big]_{i=j=0} \,\Big/~ \big[w^y_{0j0}\big]_{i=0} \,\Big/~ \big[w^x_{0i0}\big]_{j=0} \,\Big/} \\
      & \min~\langle~ & \theta_{(i-1)(j-1)}+c_{ij}, \\
      & & \min p\mapsto\theta_{pj}+w^x_{pij}, \\
      & & \min q\mapsto\theta_{iq}+w^y_{qji} ~\rangle
  \end{array}
\]
\vspace{-3mm}
\caption{\label{evaluation:gap spec}
  Specifications for the Gap problem.}
\end{figure}

\begin{paragraph}{Parenthesis problem.} Our running example; Compute
an optimal placement of parentheses in a long chain of multiplication, e.g. of matrices, where the inputs
are cost functions $x_i$ for accessing the $i$-th element and
$w_{ikj}$ for multiplying elements $[i,k)$ by elements $[k,j)$.
\end{paragraph}

\begin{paragraph}{Gap problem.}
A generalized minimal edit distance problem. Given two input strings 
$\overline{x}=x_1\cdots x_m$ and $\overline{y}=y_1\cdots y_n$,
compute the cost of transforming $x$ into $y$ by any combination of the
following steps:
  (i) Replacing $x_i$ with $y_j$, at cost $c_{ij}$,
  (ii) Deleting $x_{p+1}\cdots x_q$, at cost $w^x_{pq}$,
  (iii) Inserting $y_{p+1}\cdots y_q$ in $\overline{x}$, at cost $w^y_{pq}$.
The corresponding recurrence is shown in \Cref{evaluation:gap spec}.
\end{paragraph}

\begin{paragraph}{Protein Accordion Folding problem.} A protein can be viewed
as a string $\mathcal{P}_{1..n}$ over an alphabet of amino acids. 
The protein folds itself in a way that minimizes potential energy.
Some of the acids are {\em hydrophobic}; minimization of the total hydrophobic
area exposed to water is a major driving force of the folding process.
One possible model is packing $\mathcal{P}$ in a two-dimensional square lattice
in a way that maximizes the number of pairs of hydrophobic elements,
where the shape of the fold is an {\em accordion}, alternating between going down and going
up.
\end{paragraph}

\medskip
We also exercised our system on a number of textbook problems:
the Longest Common Subsequence (LCS) problem, the Knapsack problem,
and the Bitonic Traveling Salesman problem.
\cbdelete\diffnote{\raggedright usability estimation moved to later}

\subsection{Implementation Details}
\cbstart\diffnote{separated implementation details from experiments}
The tactic application engine is implemented in Scala. We implemented a prototype
IDE using HTML5 and AngularJS, which communicates with the engine by sending
and receiving program terms serialized as JSON. Our system supports using either
Z3 or CVC4 as the back-end SMT solver for discharging proof obligations required
for soundness proofs. Synthesis of recursive calls is done by translating the
program to \Sketch{}, which solves a correct assignment to type parameters.
To argue for the feasibility of our system, we include
SMT solver running time for the verification of the three most used tactics (figures are for CVC4),
as well as time required for \Sketch{} synthesis, in \Cref{evaluation:solving time}.
We consider an average delay of \textasciitilde 10 seconds to be reasonable, even for an interactive
environment such as Bellmania.\cbend

Tactics are implemented as small Scala classes. It is possible for
the more advanced user to extend the library by writing such classes.
To give an idea, on top of the generic TAE the {\sf Stratify} tactic was coded in 12 lines of Scala,
including the functionality that breaks a function $h$ into two functions $f$ and $g$.

The compiler back-end is implemented as another component in
\cbstart\diffnoteup{\vspace{-1cm}as part of AE, the compiler back-end was ported to Scala}Scala\cbend, processing ASTs and generating C++ code containing Intel Cilk constructs for parallelization.
It employs a thin intermediate representation layer at which the optimizations of \Cref{codegen} are applied.

\cbstart\diffnote{separated implementation details from experiments}%\vspace{-1em}% the neg vspace is only needed because of \marginpar (introduced by \diffnote)
\subsection{Experimental Results}
\label{evaluation:experimental}
\cbend

Table~\ref{evaluation:cppruntimes} shows performance improvement for our 
auto-generated implementation (AUTO) on the state-of-the-art optimized parallel
loop implementation (LOOPDP) that was used by~\cite{IPDPS15/Tithi}. It also compares AUTO with manually 
optimized recursive implementations CO\_Opt and COZ for the three problems from~\cite{IPDPS15/Tithi}\footnote{Speedups in the table are lower than in~\cite{IPDPS15/Tithi}; this should be attributed to improvement done to the
baseline LOOPDP.}. 
Our compiler automatically does \textit{copy optimization} 
as done in CO\_Opt and COZ. COZ also incorporates a low-level 
optimization of using Z-order layout of the array,
which is out of scope for this paper. %pointer arithmetics for loop traversal? 
%Explicit vectorization of loops?
$N$ is the problem size and $B$ is the base case size for using loops 
instead of recursion. It can be seen from the table that our \cbdelete\diffnote{removed remark about further optimizing by hand} implementation 
performs close to the manually optimized code.
Figure~\ref{fig:gap} depicts the performance of these implementations on one sample instance
as a function of problem size, and shows the scalability of the generated code. 

\begin{table}
\centering
\renewcommand\a{({\it i})}    % relax! it's only for this figure
\renewcommand\b{({\it ii})}
\renewcommand\c{({\it iii})}
\begin{tabular}{|l|rr@{\quad}r@{\quad}|r@{\quad}|}
  \cline{2-5}
  \multicolumn{1}{c|}{} & \multicolumn{3}{c|}{\small Verification} & \multicolumn{1}{@{\,}c@{\,}|}{\small Synthesis} \\
  \multicolumn{1}{c|}{} & \multicolumn{1}{c|}{\sf Slice} & \multicolumn{1}{c|}{\sf Stratify} & \multicolumn{1}{c|}{\sf Synth} & \multicolumn{1}{c|}{Sketch} \\
  \hline
  {\bf Paren               }  &  0.9  &  8.7   &   0.9  &  24.5     \\
  \hline
  {\bf Gap                 }  &  0.6  &  6.8   &   1.4  &  11.6     \\
  \hline
  {\bf Protein             }  &  0.9  &  3.8   &   0.7  &  9.5     \\
  \hline
  {\bf LCS                 }  &  0.9  &  1.9   &   0.5  &  3.2     \\
  \hline
  {\bf Knapsack            }  &  0.3  &  1.9   &   0.4  &  5.3     \\
  \hline
  {\bf Bitonic             }  &  0.9  &  7.2   &   0.6  &  10.1     \\
  \hline
\end{tabular}
\caption{\label{evaluation:solving time}
  Average proof search time for proof obligations and average synthesis
  time for {\sf Synth} parameters (seconds).}
\end{table}



\begin{table}
\centering
\begin{tabular}{|l|c|c|c|c|}
    \cline{2-5}
  \multicolumn{1}{c|}{} & \multicolumn{4}{c|}{\scriptsize Speedup w.r.t parallel LOOPDP on 16 cores}  \\
  \multicolumn{1}{c|}{} & \multicolumn{4}{c|}{\scriptsize   CPU (16 workers), B=64}  \\
  \multicolumn{1}{c|}{} & \multicolumn{1}{c|}{~~N~~} & \multicolumn{1}{c|}{CO\_Opt} & \multicolumn{1}{c|}{~~COZ~~} & \multicolumn{1}{c|}{AUTO}  \\
  \hline
  {\bf Parenthesis}  & 16384  & 9.8x & 11.4x & 11.1x\\
  \hline
  {\bf Gap}  & 16384 & 6.6x & 8.4x & 8.5x\\
  \hline
  {\bf Protein} & 16384  & 5.1x & 5.5x & 3.1x \\  % LOOPDP 136.579  CO 34.1431  COZ 32.1264  Bellmania (w/o manual opts) 50.986
  \hline
  {\bf LCS}  & 45000 & $-$ & $-$ & 3.9x \\
  \hline
  %{\bf Knapsack} & $-$ & $-$ & ?x\\
  %\hline
  {\bf Bitonic}  & 45000  & $-$ & $-$ & 3.8x\\
  \hline
\end{tabular}
\caption{\label{evaluation:cppruntimes}
  Performance of different C++ implementations}
\end{table}

\begin{figure}
\resizebox{8cm}{!}{
\input{gfx/charts}}
\caption{\label{fig:gap} Performance comparison for parallelized implementations for Gap problem on 16-core Intel Xeon 2.4GHz CPU}
\end{figure}

\newcommand\isum[1]{\pgfmathparse{int(#1)}\pgfmathresult}

\begin{table}
\centering
\begin{tabular}{|l|r@{\quad}r@{\quad}|r@{\quad}|}
  \cline{2-4}
  \multicolumn{1}{c|}{} & \multicolumn{2}{c|}{\small Conceptual} & \multicolumn{1}{c|}{\small Bellmania} \\
  \multicolumn{1}{c|}{} & \multicolumn{1}{c}{\# phases} & \multicolumn{1}{c|}{\# steps} & \multicolumn{1}{c|}{\# tactics} \\
  \hline
  {\bf Paren               }  &   3   &   \isum{4 + 9 + 9}   &   \isum{4 + 14 + 12}   \\
  \hline
  {\bf Gap                 }  &   3   &   \isum{9 + 9 + 9}   &   \isum{19 + 17 + 17}  \\
  \hline
  {\bf Protein             }  &   4   &   \isum{5 + 7 + 7 + 9}  &   \isum{10 + 11 + 12 + 14} \\
  \hline
  {\bf LCS                 }  &   1   &   5             &   5   \\
  \hline
  {\bf Knapsack            }  &   2   &    \isum{7 + 9} &   \isum{25 + 24}       \\
  \hline
  {\bf Bitonic             }  &   3   &    \isum{5 + 5 + 6} &   \isum{11 + 10 + 11}  \\
  \hline
\end{tabular}
\caption{\label{evaluation:user effort}
  Sizes of synthesis scripts compared to conceptual problem size (see \Cref{evaluation:effort}).}
\end{table}

\subsubsection{Estimation of User Effort}
\label{evaluation:effort}

Because Bellmania is an interactive system, we try to give a measure as to how much effort a typical user has to invest
to complete a development for the DP algorithms that comprise our test suite.
To get an idea of how domain experts think about the problem, we consult~\cite{PPoPP16/Chowdhury},
where descriptions are conveniently provided in the form of data-flow diagrams for each step of computation.
An example for such a diagram, corresponding to our \Cref{overview:recursive-A}, is shown in \Cref{evaluation:par-A-diag}; in this case, we count it as 4 steps.

%\pagebreak
%\begingroup  % want to do it but causes fig to move to the end :/
\setlength{\columnsep}{3mm}% anyway it's too late for \columnsep to affect anything at this point
\setlength{\intextsep}{3mm}
%
\begin{wrapfigure}{r}{2.8cm}
\includegraphics[width=.15\textwidth,trim=0mm 0mm 2mm 0mm]{img/par-A.pdf}
\caption[An example diagram]{\label{evaluation:par-A-diag}\\
  An example \\diagram (from~\cite{PPoPP16/Chowdhury})}
\end{wrapfigure}
%
%\endgroup

\cbstart\diffnote{\vspace{2mm}explained the numbers, and admitted the need for a more thorough study}%
We compare the sizes of these diagrams, which we label ``\# steps'', with the number of tactic
applications required to derive the respective implementation in Bellmania.
The results of the comparison are given in \Cref{evaluation:user effort}, where ``\# phases''
indicates how many recursive subroutines are included in the algorithm description 
(and in the respective development) and the two other columns give the sizes of the description vs. that of the development.
The development size is within $2\times$ of the diagrams' size in most cases, peaking at $3\times$.

Diagrams and complete Bellmania transcript for the running example are included in \Cref{annex:example:paren}.

Admittedly, the measured volume ratio provides a crude proxy for usability.
A controlled user study should be run to provide true evidence,
but that goes beyond the scope of this paper.
Still, the results are reassuring since through the course of developing the
tool we were able to reduce the number of reasoning steps from hundreds to
a few dozens, and also make them high-level.
We were also able to detect mistakes in existing (hand written) formulations
of divide-and-conquer algorithms, that eluded discovery by experts.
The additional examples in the appendix
should furnish sufficient intuition to conceive a qualitative opinion on
the efficacy of the approach and the prototype.
\cbend

%~\\~  % waste to lines to get to the end of the wrapfig
% (\vspace doesn't work because wrapfig only counts actual lines)

\input{08-further}
\section{Related Work}
\label{related}

Classical work by Smith \etal~\cite{AI85/Smith} presents rule-based transformation, stringing it
tightly with program verification. This lay the foundation for semi-automatic programming~\cite{CPS91/Blaine,TSE90/Smith,TPHOLs96/Butler}.
\cbstart\diffnote{incorporated additional references from review}
Later, these concepts have been applied to divide-and-conquer in a disciplined way
in \cite{PASCO97/Teo,EuroPar96/Chin}; these address divide-and-conquer in the
classical sense of \cite{09/CLRS} (Chapter 4), focusing on parallelism.
In Bellmania, more focus is put on re-ordering of array reads and writes, following
and mechanizing techniques related to DP from \cite{SODA06/Chowdhury,SPAA08/Chowdhury}.
In fact, traditional parallelism is taking ``for granted'' for our aggregation
operators, since they are associative and methods such as \cite{PLDI94/Fisher} apply rather trivially, and translated into the tactics \textsf{Slice}, \textsf{Assoc}, and \textsf{Distrib}.
On top of these algebraic transformations, Bellmania allows clever re-orderings,
especially through \textsf{Stratify} and \textsf{Let}.
(Some of the these tactics appear in the appendix.)
\cbend
More recently, a similar approach was introduced into Leon~\cite{OOPSLA13/Kneuss}, leveraging deductive
tools as a way to boost {\cegis}, thereby covering more programs. Bellmania takes a dual approach, where
automated techniques based on SMT are leveraged to support and improve deductive synthesis.

Inductive synthesis has been the focus of renewed interest thanks to the discovery of techniques that leverage SAT/SMT solvers to symbolically represent and search very large spaces of possible programs~\cite{APLAS09/Solar-Lezama, PLDI11/Gulwani, Onward13/Torlak},
and the use of counterexample-guided inductive synthesis ({\cegis}), which allows one to leverage inductive techniques to find programs that satisfy more general specifications. 
Our work is also inspired by the StreamBit project~\cite{PLDI05/Solar-Lezama}, which
introduced the idea of transformation rules with missing details that can be inferred by a symbolic search procedure.

Fiat~\cite{POPL15/Delaware} is another recent system that admits stepwise transformation of specifications
into programs via a refinement calculus. While Bellmania offloads proofs to SMT and \Sketch{},
Fiat uses decision procedures in Coq,
reling heavily on deductive reasoning and uses Ltac scripts for automation.
The intended users of Fiat is regular software developers who invoke pre-packaged scripts,
whereas Bellmania targets domain experts who exercise more control over the generated code.

Broadly speaking, the Bellmania system could have been implemented as a library on top of a framework
such as Coq or Why3~\cite{ESOP13/Filliatre} using binding to SMT solvers provided by these frameworks.
The decision not to do so was merely a design choice, to facilitate easier integration with our UI and with \Sketch{}.

%Polyhedral compilers offer some optimizations for the same domain of problem via tiling~\cite{HPC10/Pouchet,PLDI08/Bondhugula}.
%While showing significant speedups, these compilers cannot produce divide-and-conquer optimizations,
%which were proved to be more effective by \cite{IPDPS15/Tithi}.

Autogen~\cite{PPoPP16/Chowdhury} is a most recent advance that employs dynamic analysis to discover
a program's access pattern and learn a decomposition that can be used to generate a divide-and-conquer
implementation. The two methods are complementary,
\cbstart\diffnote{this seemed to have been unclear and it is a crucial difference between the two methods}
since Autogen does not provide correctness guarantees:
it works for a class of problems that obey a ``small world assumption'', meaning that all possible behaviors are demonstrated by traces of bounded, known size.
Crucially, in Autogen it is the user's responsibility to determine whether the input problem falls within this category;
if it does not, Autogen will not fail but instead generate an incorrect implementation.
This is a fundamental difference stemming from Autogen's use of dynamic traces \textit{vs.} purely deductive reasoning in Bellmania.
Still, the user might be able to use insights from Autogen to develop verified code in Bellmania,
where size bounds are not required.
\cbend

Pu \etal{}~\cite{OOPSLA11/Pu} have shown that recurrences for DP can be generated automatically from a non-recursive specification of the optimization problem.
This is orthogonal; in Bellmania, the recurrence is the input, and the output is an efficient divide-and-conquer implementation.
\cbstart
Obviously, the recurrence produced by \cite{OOPSLA11/Pu} can be used as input to Bellmania,
providing an even higher-level end-to-end reasoning.
\cbend

\begin{comment}
Our ``$\big/$'' operator can be compared to the separating disjunction ``$\ast$'' of Separation Logic~\cite{LICS02/Reynolds},
used to frame parts of the dynamic heap (which can be thought of as one large array),
in particular while checking that a program only accesses the parts allocated to it in its precondition.
While $\ast$ has the semantics of an existentially quantified predicate, Bellmania uses type qualifiers
to explicitly specify a formula defining each part. In this sense, it is more closely related to
Region Logic~\cite{ECOOP08/Banerjee}. These formulas make encoding in first-order logic straightforward,
and the use of Liquid Types allows for any number of dimensions and for decidable checking of domain inclusion
and disjointness.
\end{comment}
\section{Conclusion}
\label{conc}

\cbstart\diffnote{\vspace{0mm}Completely changed the narrative, and alluded to applicability in different domains as future work}%
The examples in this paper show that a few well-placed tactics can cover a wide range
of program transformations. The introduction of solver-aided tactics allowed us to make
the library of tactics smaller, by enabling the design of higher-level, more generic
tactics. Their small number gives the hope that end-users with some mathematical background
will be able to use the system without the steep learning curve that is usually associated
with proof assistants. This can be a valuable tool for algorithms research.

Moreover, limiting the number of tactics shrinks the space in which to search for programs,
so that an additional level automation may be achieved via AI or ML methods. As more
developments are done by humans and collected in a database, those algorithms would become
more adept in predicting the next step of the construction.

In a broader outlook, the technique for formalizing transformation tactics is not
particularly attached to divide-and-conquer algorithms and their implementations.
In this work, we constructed a generic tactic application engine, on top of which
the tactics required for our domain were easy to implement.
This gives rise to the hope that, in the future, the same approach can be applied
to other domains, in the interest of encoding domain knowledge, providing better
DSLs that offer users the power to write high-level programs without sacrificing
performance.
\cbend

\section*{Acknowledgments}

This work is supported by NSF Grants CCF-1139056, CCF-1439084 and CNS-1553510. We thank Shoaib Kamil and Adobe research for their feedback and support.


\bibliographystyle{abbrvnat}
\bibliography{oopsla2016}
% The bibliography should be embedded for final submission.

%\begin{thebibliography}{}
%\softraggedright

%\bibitem[Smith et~al.(2009)Smith, Jones]{smith02}
%P. Q. Smith, and X. Y. Jones. ...reference text...

%\end{thebibliography}

% This goes to supplementary material
\clearpage

\appendix

\input{A-proof}
\vfill\pagebreak
\section{More Tactics}
\label{annex:more tactics}

\newcommand\self{$\star$}

For most of the tactics below, the proof obligation is exactly the equality that that
expresses the rewrite. We signify this with the notation \Obligations \self.

\tacticdef{Shrink} \label{tactics:Shrink}
\[f ~=~ f :: \T\] 

Used to specify tighter qualifiers for the type of a sub-term.

\Obligations \self.

For arrow-typed terms, this essentially requires to prove that $f$ 
is only defined for arguments in the domain of $\T$, and that the values are in the range of $\T$.
This can be seen as a special case of {\sf Slice} with $r=1$, with the additional
feature of specifying the range as well.

\tacticdef{Associativity}

\[\reduce\big\langle \reduce\langle\overline x_1\rangle,\cdots, \reduce\langle\overline x_r\rangle\big\rangle ~=~ \reduce \langle \overline x_1, \cdots, \overline x_r\rangle\]
%
where $\reduce$ is a built-in aggregation ($\min$, $\max$, $\Sigma$), 
and $\overline x_i$ are lists of terms (of the same type).
If any of $\overline x_i$ is of length one, $\reduce\langle\overline x_i\rangle$ can be
replaced by $\overline x_i$.

\Obligations none.

\tacticdef{Distributivity}\vspace{3mm}
Let $e$ be an expression with a hole, $e[\square] = (\cdots \square \cdots)$.
%
\[\renewcommand\arraystretch{1.2}
  \begin{array}{@{}r@{~}c@{~}l@{}}
    e[t_1/\cdots/t_r] &=& e[t_1] / \cdots / e[t_r] \\
    e[t_1/\cdots/t_r] &=& \reduce\langle e[t_1],\cdots,e[t_r]\rangle \\
    \reduce e[t_1/\cdots/t_r] &=& \reduce\langle \reduce e[t_1],\cdots,\reduce e[t_r]\rangle
  \end{array}\]

This tactic provides several alternatives for different uses of aggregations.
Clearly, $\big/$ does not distribute over any expression; we give just a few examples
where this tactic is applicable.
\begin{itemize}
  \item $(x/y)+1 ~=~ (x+1)~/~(y+1)$
  \item $x/0 ~=~ \max\langle x,0\rangle$ ~(for $x:\N$)
  \item $\min \big([f]_{J_0}~\big/~[f]_{J_1}\big) ~=~
         \min\left\langle \min [f]_{J_0} ~,~ \min [f]_{J_1}\right\rangle$
\end{itemize}

\Obligations \self.

\tacticdef{Elimination}
\[e[t] ~=~ e[\bot]\]
%
Used to eliminate a sub-term that is either always undefined or has no effect
in the context in which it occurs.

\Obligations \self.

\tacticdef{Let Insertion}\vspace{3mm}
Let $e$ be an expression with a hole, $e[\square] = (\cdots x_1 \mapsto \cdots x_k\mapsto \cdots \square \cdots)$, 
where $x_{1..k}\mapsto$ are abstraction terms enclosing $\square$. The bodies may contain arbitrary terms
in addition to these abstractions.
%
\[\renewcommand\arraystretch{1.2}
  \begin{array}{@{}r@{}c@{}l@{}}
  e[t] & ~=~ & (\overline{x}\mapsto t) ~\applt~ z\mapsto e[z\,\overline{x}] \\
  e[\reduce \langle\overline{a},\overline{b}\rangle] & ~=~ &
     (\overline{x}\mapsto \reduce\langle\overline{a}\rangle) \\ 
     ~&~&~~~~~~~\applt~ z\mapsto e[\reduce\langle z\,\overline{x}, \overline{b}\rangle]
  \end{array}\]
%
where $\overline{x}=x_{1..k}$, and $z$ is a fresh variable.
This tactic also has a special version that involves $\reduce$.
The items in $\langle \overline{a},\overline{b}\rangle$ may be interleaved, since
$\min,\max,\Sigma$ all happen to be commutative.\footnote{If non-commutative functions get added in the future, then this will change into $\langle \overline{a}, \overline{b}, \overline{c}\rangle$ non-interleaving, with the right hand side being $(\overline{x}\mapsto \reduce\langle\overline{b}\rangle) ~\applt~ z\mapsto e[\reduce\langle \overline{a}, z\,\overline{x}, \overline{c}\rangle]$.}

\Obligations tactic, if $z$ occurs free in $e$; otherwise none.

\tacticdef{Let Insertion [$\reduce$]}
%
%\[e[\reduce \langle\overline{a},\overline{b}\rangle] ~=~ 
%  (\overline{x}\mapsto \reduce\langle\overline{a}\rangle) ~\applt~ z\mapsto e[\reduce\langle z\,\overline{x}, \overline{b}\rangle]
%  \]
%
\[\renewcommand\arraystretch{1.2}
  \begin{array}{@{}r@{}c@{}l@{}}
  e[\reduce \langle\overline{a},\overline{b}\rangle] & ~=~ &
     (\overline{x}\mapsto \reduce\langle\overline{a}\rangle) \\ 
     ~&~&~~~~~~~\applt~ z\mapsto e[\reduce\langle z\,\overline{x}, \overline{b}\rangle]
  \end{array}\]
%

where $\overline{x}=x_{1..k}$, and $z$ a fresh variable.

\Obligations \self, if $z$ occurs free in $e$; otherwise none.

\tacticdef{Padding}
\[t ~=~ \big(t~/~f_1/\cdots/f_r\big) :: \T\]
%
where $\T$ is the type of $t$. This tactic is commonly used with Let insertion,
to make the type of a sub-computation match the type of the entire term.

\Obligations \self.

\tacticdef{Pull Out}\vspace{3mm}
For $e[\square]$ as defined previously:
\[z ~=~ \overline{x}\mapsto t\]
%
where $z$ is a fresh variable.

Similar to Let Insertion, but does not change the original term; instead,
it is used to single out and name a particular expression $t$, preserving the context
in which it occurs in $e[t]$. It is not a tactic {\it per se}, as it does not
actually effect any transformation on $e[t]$; instead, it is designed to increase
readability of the development and simplify successive human-computer interaction.

\vfill\pagebreak
\section{Examples}

Many aspects of Bellmania are best illustrated via examples.
While the main sections include many such examples,
some more elaborated ones may prove useful and interesting.

\subsection{Development of Parenthesis (with diagrams)}
\label{annex:example:paren}

The running example from \Cref{intro,overview} has a total of three subroutines, which we label A, B, and C.
As promised, we include original design diagrams by the technique's authors taken from~\cite{PPoPP16/Chowdhury}.
The blocks (triangles) in the diagrams represent intermediate steps of the computation.
Below each diagram, we show a transcript of what the user has to type in when using Bellmania
to carry out the same development.

Boxed letters in the scripts are used to refer to sub-terms of the current program
(in \Cref{tactics} we used boxed digits, but in the actual UI we use letters
because there are more of them).
Reading the scripts in a non-interactive setting might be hard since the reader cannot
observe the program; they are listed here just to give an idea of the size and structure.
Bellmania repository~\cite{bellmaniarepo} contains some screen-shots of an interactive session.

\newcommand\xw[1]{w^x_{#1}}
\newcommand\yw[1]{w^y_{#1}}

\begin{figure}
\begin{center}
\includegraphics[width=.15\textwidth]{img/par-A.pdf}
\end{center}
%
\begin{lstlisting}[language=bellmania]
Slice (find (\theta |-> ?)) (? <|J_0*J_0, J_0*J_1, J_1*J_1|>)
Stratify "/" (fixee [A]) [A] \psi
Stratify "/" (fixee [A]) [A] \psi
[A] [B] [C] |-> SynthAuto . ... \psi
\end{lstlisting}
\caption{\label{annex:Paren-A}
Development of subroutine A of the Parenthesis problem as conceptually described in~\cite{PPoPP16/Chowdhury} (top) and using Bellmania (bottom).}
\end{figure}

\begin{figure}
\begin{center}
\includegraphics[width=.3\textwidth]{img/par-B.pdf}
\end{center}
%
\begin{lstlisting}[language=bellmania]
Slice f (? <|K_0*K_2, K_0*K_3, K_1*K_2, K_1*K_3|>)
[D] |-> Stratify "/" (fixee .) . \psi
[C] |-> Stratify "/" (fixee .) . \psi
[E] |-> Stratify "/" (fixee .) . \psi

<|Slice (find  (k |-> ?)) <|K_0, K_1, K_2, K_3|>,
 Slice (find  (k |-> ?)) <|K_1, K_2, K_3|>,
 Slice (find  (k |-> ?)) <|K_0, K_1, K_2|> |>

Distrib min
Assoc min

<|Stratify min (fixee $\lbox A$) <|$\lbox G$, $\lbox J$|> \psi,
 Stratify min (fixee $\lbox B$) <|$\lbox M$, $\lbox O$|> \psi,
 Stratify min (fixee $\lbox C$) <|$\lbox R$, $\lbox T$|> \psi |>
Stratify min (fixee $\lbox A$) $\langle \lbox{\,I\,}, \lbox K\rangle$ $\psi$

$\lbox{\,I\,}$ $\lbox S$ $\lbox Z$ $\underline{\lbox G}$ $\lbox M$ $\lbox P$ $\lbox W$ $\underline{\lbox D}$ |-> SynthAuto . ... \psi
\end{lstlisting}
\caption{\label{annex:Paren-B}
Same, for subroutine B of Parenthesis.}
\end{figure}

\begin{figure}
\begin{center}
\includegraphics[width=.3\textwidth]{img/par-B.pdf}
\end{center}
%
\begin{lstlisting}[language=bellmania]
Slice (find (i |-> ?)) <|L_0*L_4,L_0*L_5,L_1*L_4,L_1*L_5|>
Let "/" (slasher $\lbox A$) $\lbox A$ $\psi$
Let "/" (slasher $\lbox A$) $\lbox A$ $\psi$
Let "/" (slasher $\lbox A$) $\lbox A$ $\psi$
Slice (findAll ($k \mapsto {?}$)) $\langle L_2, L_3\rangle$

Distrib min
Assoc min

<|Let min (slasher $\lbox A$) $\langle\lbox E, \lbox G\rangle$ $\psi$,
 Let min (slasher $\lbox B$) $\langle\lbox H,\lbox {\,J\,}\rangle$ $\psi$,
 Let min (slasher $\lbox C$) $\langle\lbox K,\lbox M\rangle$ $\psi$,
 Let min (slasher $\lbox D$) $\langle\lbox N,\lbox P\rangle$ $\psi$ $\rangle$

$\lbox A$ $\lbox B$ $\lbox C$ $\lbox D$ $\lbox E$ $\lbox F$ $\lbox G$ $\lbox H$ |-> SynthAuto . ... $\psi$
\end{lstlisting}
\caption{Same, for subroutine C of Parenthesis}
\end{figure}

\vfill

\subsection{Qualified Type Inference}
\label{annex:example:type inference}

We provide an example of how qualifiers are inferred in program terms.

\medskip
\exampleTitle

\noindent Assume that:
\begin{itemize}
 \item $I$, $T$ are types
 \item $\widehat{I}_0 : I\to\B$ is a unary qualifier
 \item $0:T$ is a constant
 \item $\S$ a type variable, 
\end{itemize}
 
Consider the term $(f:I_0\to\S)\,i\mapsto f\,i\,i ~/~ 0$.
The first step of Hindley-Milner inference will induce the following type shapes through unification:

\begin{tikzpicture}[node distance=1em]
  \node(farg) {$(f:I\to\S)$};
  \node(iarg)[right=of farg] {$i$};
  \node(mapsto)[right=of iarg] {$\mapsto$};
  \node(f)[right=of mapsto] {$f$};
  \node(i1)[right=of f] {$i$};
  \node(i2)[right=of i1] {$i$};
  \node(slash)[right=of i2] {$\big/$};
  \node(zero)[right=of slash] {$0$};
  \node(l0)[coordinate,below=1mm of farg] {};
  \node(l1)[coordinate,below=4.5mm of l0] {};
  \node(l2)[coordinate,below=3.5mm of l1] {};
  \node(l3)[coordinate,below=3.5mm of l2] {};
  \node(l4)[coordinate,below=3.5mm of l3] {};
  \def\ytip{2pt}
  \draw (l0 -| farg.west) -- +(0,-\ytip) -| (l0 -| farg.east) node[pos=0.25,below] {\tiny $I\to I\to T$};
  \draw (l0 -| iarg.west) -- +(0,-\ytip) -| (l0 -| iarg.east) node[pos=0.25,below] {\tiny $I$};
  \draw (l0 -| f.west) -- +(0,-\ytip) -| (l0 -| f.east) node[pos=0.5,anchor=23,inner xsep=0] {\tiny $I\to I\to T$};
  \draw (l0 -| i1.west) -- +(0,-\ytip) -| (l0 -| i1.east) node[pos=0.25,below] {\tiny $I$};
  \draw (l0 -| i2.west) -- +(0,-\ytip) -| (l0 -| i2.east) node[pos=0.25,below] {\tiny $I$};
  \draw (l0 -| zero.west) -- +(0,-\ytip) -| (l0 -| zero.east) node[pos=0.25,below] {\tiny $T$};
  \draw (l1 -| f.west) -- +(0,-\ytip) -| (l1 -| i1.east) node[pos=0.25,below] {\tiny $I\to T$};
  \draw (l2 -| f.west) -- +(0,-\ytip) -| (l2 -| i2.east) node[pos=0.25,below] {\tiny $T$};
  \draw (l3 -| f.west) -- +(0,-\ytip) -| (l3 -| zero.east) node[pos=0.25,below] {\tiny $T$};
  \draw (l4 -| farg.west) -- +(0,-\ytip) -| (l4 -| zero.east) node[pos=0.25,below] {\tiny $(I\to I\to T)\to I\to T$};
  
  % numerals!
  \node[yshift=-1pt] at (farg.north east) {\tiny 0};
  \node at (iarg.north east) {\tiny 0};
  \node at (i1.north east) {\tiny 1};
  \node at (i2.north east) {\tiny 2};
  \node at (f.north east) {\tiny 1};
\end{tikzpicture}

Superscript numerals denote different occurrences of the same variable.
In this case, the type variable $\S$ has been assigned $I\to T$.

The process would have stopped here if it weren't for the qualifier $I_0$ used in the type for $f$.
At this point we can use type refinements to get more accurate types for $f$ and $i$
in the body of the function term.
%
\newcommand\oops[1]{^{\,^{\mbox{\tiny #1}}}}
\[
\frac{f:I_0\to I\to T, ~ i:I ~~\vdash~~ f\oops1 : I\to I\to T}
     {f:I_0\to I\to T, ~ i:I ~~\vdash~~ f\oops1 : I_0\to I\to T}
 \]
 
Notice that $(I\to I\to T)\sqcap (I_0\to I\to T) = I_0\to I \to T$. Truthfully, in this case this is quite a trivial result.

\smallskip
Let $\Gamma = \{f:I_0\to I\to T, ~ i:I\}$.
%
\[
\frac{\begin{array}{rcl} \Gamma & \vdash & (f\oops1\,i\oops1) : I\to T, 
                                        ~~ f\oops1: I_0\to I\to T,      
                                        ~~ i\oops1 : I                 \end{array}}
     {\begin{array}{rcl} \Gamma & \vdash & (f\oops1\,i\oops1) : I\to T, 
                                        ~~ f\oops1: I_0\to I\to T,      
                                        ~~ i\oops1 : I_0
     \end{array} }
     \]

The types of $f\oops1$ and $f\oops1\,i\oops1$ have not changed, but the type of $i\oops1$ was lowered to $I_0$.

\medskip
After applying the typing rules similarly to all the sub-terms, we get the inferred types as shown:

\begin{tikzpicture}[node distance=1em]
  \node(farg) {$(f:I_0\to\S)$};
  \node(iarg)[right=of farg] {$i$};
  \node(mapsto)[right=of iarg] {$\mapsto$};
  \node(f)[right=of mapsto] {$f$};
  \node(i1)[right=of f] {$i$};
  \node(i2)[right=of i1] {$i$};
  \node(slash)[right=of i2] {$\big/$};
  \node(zero)[right=of slash] {$0$};
  \node(l0)[coordinate,below=1mm of farg] {};
  \node(l1)[coordinate,below=4.5mm of l0] {};
  \node(l2)[coordinate,below=3.5mm of l1] {};
  \node(l3)[coordinate,below=3.5mm of l2] {};
  \node(l4)[coordinate,below=3.5mm of l3] {};
  \def\ytip{2pt}
  \draw (l0 -| farg.west) -- +(0,-\ytip) -| (l0 -| farg.east) node[pos=0.25,below] {\tiny $I_0\to I\to T$};
  \draw (l0 -| iarg.west) -- +(0,-\ytip) -| (l0 -| iarg.east) node[pos=0.25,below] {\tiny $I$};
  \draw (l0 -| f.west) -- +(0,-\ytip) -| (l0 -| f.east) node[pos=0.5,anchor=23,inner xsep=0] {\tiny $I_0\to I\to T$};
  \draw (l0 -| i1.west) -- +(0,-\ytip) -| (l0 -| i1.east) node[pos=0.25,below] {\tiny $I_0$};
  \draw (l0 -| i2.west) -- +(0,-\ytip) -| (l0 -| i2.east) node[pos=0.25,below] {\tiny $I$};
  \draw (l0 -| zero.west) -- +(0,-\ytip) -| (l0 -| zero.east) node[pos=0.25,below] {\tiny $T$};
  \draw (l1 -| f.west) -- +(0,-\ytip) -| (l1 -| i1.east) node[pos=0.25,below] {\tiny $I_0\to T$};
  \draw (l2 -| f.west) -- +(0,-\ytip) -| (l2 -| i2.east) node[pos=0.25,below] {\tiny $T$};
  \draw (l3 -| f.west) -- +(0,-\ytip) -| (l3 -| zero.east) node[pos=0.25,below] {\tiny $T$};
  \draw (l4 -| farg.west) -- +(0,-\ytip) -| (l4 -| zero.east) node[pos=0.25,below] {\tiny $(I_0\to I\to T)\to I\to T$};
\end{tikzpicture}

\medskip
\hrule
\vfill

\iftr % The Simplified Arbiter example overflows the page limit (20) and will appear in the TR
\pagebreak
\subsection{Simplified Arbiter}

To give another example of a development for a different problem, we provide a simplified version of the Gap problem (\Cref{evaluation}), which we call the Simplified Arbiter.
Two processes $x$ and $y$ must be scheduled to run $n$ and $m$ seconds,
respectively, on a single processor, using one-second slots.
Execution starts at $t=0$. The cost for scheduling the slots $[a..b)$ of $x$ after
having scheduled slots $[0..c)$ of $y$
is given by $\xw{abc}$, and the cost for schedulting the slots $[a..b)$ of $y$
after scheduling $[0..c)$ of $x$ is given by $\yw{abc}$.

In Bellmania language, the Simplified Arbiter example is specified by
a base case $\Psi$ and a computation
part $A^{IJ}$, where $I$ and $J$ are row and column index sets, respectively.

\[
  \renewcommand\arraystretch{1.5}
  \begin{array}{@{}l@{}l@{}c@{}c@{}l@{}l@{}}
    \Psi ~=~ & \lspan5{\fix \Big(\theta\,i\,j\mapsto
	              [0]_{i=j=0} ~\big/~ [\yw{0j0}]_{i=0} ~\big/~ [\xw{0i0}]_{j=0}\Big)} \\
    A^{^{IJ}} ~=~ 
	      & \psi\mapsto \fix \theta\,& i & j & \mapsto\min\,\langle~ & \psi_{ij} \\
	      & & ^{^{(I)}} & ^{^{(J)}} & & \min \vtyped p I \mapsto\theta_{pj}+w_{pij}, \\
	      & & & & & \min \vtyped q J \mapsto\theta_{iq}+w'_{qji} ~\rangle
  \end{array}
\]

Vertical typeset was used to save some horizontal space, but $\vtyped v\T$
should be read as just $v:\T$.

\medskip

After Richard applies {\sf Slice}, he gets the four quadrants $I_0\times J_0$, $I_0\times J_1$, $I_1\times J_0$, $I_1\times J_1$
(\Cref{overview:quadrants}). The system defined unary qualifiers with the axioms:
\[
\begin{array}{c@{\qquad}c}
  \forall i{:}I.~~I_0(i)\lor I_1(i)   &    \forall i_0{:}I_0,~i_1{:}I_1.~~i_0<i_1 \\
  \forall j{:}J.~~J_0(j)\lor J_1(j)   &    \forall j_0{:}J_0,~j_1{:}J_1.~~j_0<j_1 \\
\end{array}
\]

The program is just about to
grow quite large; to make such terms easy to read and refer to, we provide
boxed letters as labels for sub-terms, using them as abbreviations where they
occur in the larger expression

In addition, to allude to the reader's intuition, expressions of the form
$a/b/c/d$ will be written as $\quadrants*{a}{b}{c}{d}$ when the slices
represent quadrants.

\begin{tacticbox}{Slice}
  \begin{array}{@{} l  l @{}}
       f ~=~ \theta\,i\,j\mapsto \cdots \\
       X_1 ~=~ \_\times I_0\times J_0 &
       X_2 ~=~ \_\times I_0\times J_1 \\
       X_3 ~=~ \_\times I_1\times J_0 &
       X_4 ~=~ \_\times I_1\times J_1 \\[.5em]
       \cspan2{\mbox{\small ({\it each} ``\_'' {\it is a fresh type variable})}}
  \end{array}
\end{tacticbox}

\begin{equation}
  \renewcommand\arraystretch{1.2}
  \begin{array}{@{}r@{}l@{}c@{}c@{}l@{}l@{}}
    A^{^{IJ}} =~ & \lspan5{\psi\mapsto \fix \quadrants{\lbox A}{\lbox B}{\lbox C}{\lbox D}} \\
	\lbox A ~=~ & \theta\,& i & j & \mapsto\min\,\langle~ & \psi_{ij} \\
	      & & ^{^{(I_0)}} & ^{^{(J_0)}} & & \min \vtyped p I \mapsto\theta_{pj}+w_{pij}, \\
	      & & & & & \min \vtyped q J \mapsto\theta_{iq}+w'_{qji} ~\rangle \\
	\lbox B ~=~ & \theta\,& i & j & \mapsto\min\,\langle~ & \psi_{ij} \\
	      & & ^{^{(I_0)}} & ^{^{(J_1)}} & & \min \vtyped p I \mapsto\theta_{pj}+w_{pij}, \\
	      & & & & & \min \vtyped q J \mapsto\theta_{iq}+w'_{qji} ~\rangle \\
	\lbox C ~=~ & \theta\,& i & j & \mapsto\min\,\langle~ & \psi_{ij} \\
	      & & ^{^{(I_1)}} & ^{^{(J_0)}} & & \min \vtyped p I \mapsto\theta_{pj}+w_{pij}, \\
	      & & & & & \min \vtyped q J \mapsto\theta_{iq}+w'_{qji} ~\rangle \\
	\lbox D ~=~ & \theta\,& i & j & \mapsto\min\,\langle~ & \psi_{ij} \\
	      & & ^{^{(I_1)}} & ^{^{(J_1)}} & & \min \vtyped p I \mapsto\theta_{pj}+w_{pij}, \\
	      & & & & & \min \vtyped q J \mapsto\theta_{iq}+w'_{qji} ~\rangle \\
  \end{array}
  \label{evaluation:A sliced}
\end{equation}

\begin{tacticbox}{Let}
   e[\square] ~=~ \quadrants*{\square}{\lbox*B}{\lbox*C}{\lbox*D} \qquad
   t ~=~ \lbox A
\end{tacticbox}

\begin{equation}
  A^{^{IJ}} =~ \psi\mapsto \fix \left(\lbox A \applt z\mapsto\quadrants{z}{\lbox B}{\lbox C}{\lbox D}\right)
\end{equation}

\begin{tacticbox}{Stratify[with Padding]}
  \begin{array}{@{} l @{} l @{}}
    f ~=~ \quadrants*{\lbox*A}{\lifted\psi}{\lifted\psi}{\lifted\psi}
         & \mbox{\small ({\it recall that } $\lifted\psi=\theta\mapsto\psi$)} \\
    g ~=~ z\mapsto\quadrants*{z}{\lbox*{B}}{\lbox*C}{\lbox*D} &
    \qquad\psi=\psi
  \end{array}
\end{tacticbox}

\begin{equation}
  A^{^{IJ}} =~ \psi\mapsto \fix \quadrants{\lbox A}{\lifted\psi}{\lifted\psi}{\lifted\psi} ~\applt~ \psi\mapsto\fix\quadrants{\lifted\psi}{\lbox B}{\lbox C}{\lbox D}
  \label{evaluation:stratify A 1}
\end{equation}

Notice that an existing variable $\psi$ is reused, rebinding any occurrences within $\lbox B$, $\lbox C$, $\lbox D$.
This effect is useful, as it limits the context of the expression: the inner $\psi$ shadows the outer $\psi$,
meaning $\lbox B$, $\lbox C$, $\lbox D$ do not need to access the data that was input to $\lbox A$, only its
output.

\medskip
The sequence Let, Stratify[with Padding] is now applied in the same manner to $\lbox B$
and $\lbox C$ (see \Cref{overview:slice-stratify-synth}). We do not list the applications as they are analogous to the previous ones.

\begin{equation}
  \renewcommand\arraystretch{1.5}
  \begin{array}{l@{}l}
    A^{^{IJ}} =~ \psi\mapsto{} & \fix \quadrants{\lbox A}{\lifted\psi}{\lifted\psi}{\lifted\psi} ~\applt~ 
                 \psi\mapsto\fix\quadrants{\lifted\psi}{\lbox B}{\lifted\psi}{\lifted\psi} ~\applt \\
               & \psi\mapsto\fix\quadrants{\lifted\psi}{\lifted\psi}{\lbox C}{\lifted\psi} ~\applt~
                 \psi\mapsto\fix\quadrants{\lifted\psi}{\lifted\psi}{\lifted\psi}{\lbox D}
  \end{array}
\end{equation}

\begin{tacticbox}{Synth}
	\begin{array}{@{}l@{}c@{}c@{}l@{}l}
       \lspan5{h_1= \lbox A} \\
       \lspan5{h_{2,3,4}=\lifted\psi} \\
	   f_1 = \theta\,& i & j & \mapsto\min\,\langle~ & \psi_{ij} \\
	      & ^{^{(I_0)}} & ^{^{(J_0)}} & & \min \vtyped p {I_0} \mapsto\theta_{pj}+w_{pij}, \\
	      & & & & \min \vtyped q {J_0} \mapsto\theta_{iq}+w'_{qji} ~\rangle \\
	   \lspan5{f_{2,3,4} = \lifted\psi}
   \end{array}
\end{tacticbox}

\begin{equation}
  \renewcommand\arraystretch{1.5}
  \begin{array}{l@{}l}
    A^{^{IJ}} =~ \psi\mapsto{} & \quadrants{A^{^{I_0J_0}}_\psi\!\!\!}{\psi}{\psi}{\psi} ~\applt~ 
                 \psi\mapsto\fix\quadrants{\lifted\psi}{\lbox B}{\lifted\psi}{\lifted\psi} ~\applt \\
               & \psi\mapsto\fix\quadrants{\lifted\psi}{\lifted\psi}{\lbox C}{\lifted\psi} ~\applt~
                 \psi\mapsto\fix\quadrants{\lifted\psi}{\lifted\psi}{\lifted\psi}{\lbox D}
  \end{array}
  \label{fix A}
\end{equation}

We note that $\fix f_1=A^{^{I_0J_0}}_\psi$ are identical (up to $\beta$-reduction),
which is the whole reason $f_1$ was chosen. Also, we took the liberty
to simplify $\fix\lifted\psi$ into $\psi$ --- although this is not necessary --- just to display
a shorter term.

\medskip
The next few tactics will focus on the subterm $\lbox B$ from \eqref{evaluation:A sliced}.

\begin{equation}
  \renewcommand\arraystretch{1.2}
  \begin{array}{@{}r@{}l@{}c@{}c@{}l@{}l@{}}
	\lbox B ~=~ & \theta\,& i & j & \mapsto\min\,\langle~ & \psi_{ij} \\
	      & & ^{^{(I_0)}} & ^{^{(J_1)}} & & \min \vtyped p I \mapsto\theta_{pj}+w_{pij}, \\
	      & & & & & \min \vtyped q J \mapsto\theta_{iq}+w'_{qji} ~\rangle
  \end{array}
\end{equation}

\begin{tacticbox}{Slice}
  \begin{array}{@{} l @{}}
    f = \vtyped q J \mapsto \theta_{iq}+w'_{qji} \\
    X_1 = J_0\to\_ \qquad X_2 = J_1\to\_
  \end{array}
\end{tacticbox}

\begin{equation}
  \renewcommand\arraystretch{1.2}
  \begin{array}{@{}r@{}l@{}c@{}c@{}l@{}l@{}l@{}}
	\lbox B ~=~ & \theta\,& i & j & \mapsto\min\,\big\langle~ & \psi_{ij} \\
	      & & ^{^{(I_0)}} & ^{^{(J_1)}} & & \lspan2{\min \vtyped p I \mapsto\theta_{pj}+w_{pij},} \\
	      & & & & & \min \big( & (\vtyped q {J_0} \mapsto\theta_{iq}+w'_{qji}) ~\big/~ \\
	      & & & & & & (\vtyped q {J_1} \mapsto\theta_{iq}+w'_{qji})\big)  ~\big\rangle
  \end{array}
  \label{evaluation:horiz sliced}
\end{equation}

For the intuition behind this, see the top-right part of \Cref{evaluation:slicing strategy}.
The colors represent cell ranges that will be read by different sub-routines (presumably running on different cores).
The range of $q$ is split into the part that lies within \qbox1 ($q\in J_0$) and the one that
lies within \lbox2 ($q\in J_1$). The same reasoning is applied to the other quadrants.

\begin{tacticbox}{Distributivity}
  \begin{array}{@{} l @{}}
    e[\square] = \min\square \\
    t_1 = \min \vtyped q {J_0} \mapsto \theta_{iq}+w'_{qji} \\
    t_2 = \min \vtyped q {J_1} \mapsto \theta_{iq}+w'_{qji} \\
  \end{array}
\end{tacticbox}

\begin{tacticbox}{Associativity}
  \begin{array}{@{} l @{} l @{}}
    \lspan2{\reduce = \min} \\
    \overline x_1 ={} & \psi_{ij} \\
    \overline x_2 ={} & \min \vtyped p I \mapsto\theta_{pj}+w_{pij} \\
    \overline x_3 ={} & \min \vtyped q {J_0} \mapsto \theta_{iq}+w'_{qji} ~, \\
                      & \min \vtyped q {J_1} \mapsto \theta_{iq}+w'_{qji}
  \end{array}
\end{tacticbox}

\begin{equation}
  \renewcommand\arraystretch{1.2}
  \begin{array}{@{}r@{}l@{}c@{}c@{}l@{}l@{}}
	\lbox B ~=~ & \theta\,& i & j & \mapsto\min\,\big\langle~ & \psi_{ij} \\
	      & & ^{^{(I_0)}} & ^{^{(J_1)}} & & \min \vtyped p I \mapsto\theta_{pj}+w_{pij}, \\
	      & & & & & \min \vtyped q {J_0} \mapsto\theta_{iq}+w'_{qji}, \\
	      & & & & & \min \vtyped q {J_1} \mapsto\theta_{iq}+w'_{qji} ~\big\rangle
  \end{array}
\end{equation}

\begin{tacticbox}{Let[$\reduce$]}
  \begin{array}{@{} r @{} l @{}}
    e[\square] ={} & \quadrants{\lifted\psi}{\theta\,i\,j\mapsto\square}{\lifted\psi}{\lifted\psi} \\
    \overline a ={} & \psi_{ij}, ~\min \vtyped q{J_0}\mapsto \theta_{iq}+w'_{qji} \\
    \overline b ={} & \min \vtyped p I \mapsto\theta_{pj}+w_{pij}, \\
                    & \min \vtyped q {J_1} \mapsto\theta_{iq}+w'_{qji}
  \end{array}
\end{tacticbox}

\begin{equation}
  \renewcommand\arraystretch{1.2}
  \begin{array}{r @{} l @{} c @{} c @{} l @{} l}
    \lspan6{
    \quadrants{\lifted\psi}{\lbox B}{\lifted\psi}{\lifted\psi} =
      \fix\left(\lbox E \applt z\mapsto\quadrants{\lifted\psi}{\lbox F}{\lifted\psi}{\lifted\psi}\right) 
    } \\[1.5em]
    ~\lbox E ={} &
      \theta & i & j & \mapsto\min\langle & \psi_{ij}, \\
             & & ^{^{(I_0)}} & ^{^{(J_1)}} &
                                          & \min \vtyped q{J_0}\mapsto\theta_{iq}+w'_{qji} \rangle \\
    ~\lbox F ={} &
      \theta & i & j & \mapsto\min\langle & z_{\theta ij}, \\
             & & ^{^{(I_0)}} & ^{^{(J_1)}} &
                                         & \min \vtyped p I \mapsto\theta_{pj}+w_{pij}, \\
             & & & &                     & \min \vtyped q {J_1} \mapsto\theta_{iq}+w'_{qji}\rangle
  \end{array}
\end{equation}


\begin{figure}
\input{gfx/evaluation-slicing-strategy}
\medskip
\caption{\label{evaluation:slicing strategy}
  The strategy for applications of {\sf Slice} in the case study.}
\end{figure}

\begin{tacticbox}{Stratify[with Padding]}
  \begin{array}{@{} l @{}}
    f ~=~ \quadrants*{\lifted\psi}{\lbox*E}{\lifted\psi}{\lifted\psi} \\
    g ~=~ z\mapsto\quadrants*{\lifted\psi}{\lbox*F}{\lifted\psi}{\lifted\psi}
    \qquad\psi=\psi
  \end{array}
\end{tacticbox}

\begin{equation}
  \renewcommand\arraystretch{1.2}
  \begin{array}{r @{} l @{} c @{} c @{} l @{} l}
    \lspan6{
    \fix\quadrants{\lifted\psi}{\lbox B}{\lifted\psi}{\lifted\psi} =
      \fix\quadrants{\lifted\psi}{\lbox E}{\lifted\psi}{\lifted\psi} \applt
      \psi\mapsto\fix\quadrants{\lifted\psi}{\lbox F}{\lifted\psi}{\lifted\psi}
    } \\[1.5em]
    ~\lbox E ={} &
      \theta & i & j & \mapsto\min\langle & \psi_{ij}, \\
             & & ^{^{(I_0)}} & ^{^{(J_1)}} &
                                          & \min \vtyped q{J_0}\mapsto\theta_{iq}+w'_{qji} \rangle \\
    ~\lbox F ={} &
      \theta & i & j & \mapsto\min\langle & \psi_{ij}, \\
             & & ^{^{(I_0)}} & ^{^{(J_1)}} &
                                          & \min \vtyped p I \mapsto\theta_{pj}+w_{pij}, \\
             & & & &                      & \min \vtyped q {J_1} \mapsto\theta_{iq}+w'_{qji}\rangle
  \end{array}
\end{equation}

\noindent
Define
\begin{equation}
  \renewcommand\arraystretch{1.2}
  \begin{array}{@{}l @{} l @{\!} c @{} c @{} l @{} l@{}}
  B^{^{IJ_0J_1}} =~ \big( & \psi\mapsto \\
      & \fix
      \theta & i & j & \mapsto\min\langle & \psi_{ij}, \\
           & & ^{^{(I)}} & ^{^{(J)}} &
                                          & \min \vtyped q{J_0}\mapsto\theta_{iq}+w'_{qji} \rangle\big) \\
      & \lspan5{:: \big((I\times J_0)\to\R\big) \to \big((I\times J_1)\to\R\big)}
  \end{array}
\end{equation}

\begin{tacticbox}{Synth}
  \begin{array}{@{} l @{} c @{} c @{} l @{} l @{}}
    \lspan5{h_2 = \lbox E} \\
    \lspan5{h_{1,3,4} = \lifted\psi} \\
    f_2 = 
      \theta & i & j & \mapsto\min\langle & \psi_{ij}, \\
             & ^{^{(I_0)}} & ^{^{(J_1)}} &
                                          & \min \vtyped q {J_0} \mapsto\theta_{iq}+w'_{qji}\rangle \\
    \lspan5{f_{1,3,4} = \lifted\psi}
  \end{array}
\end{tacticbox}

\begin{tacticbox}{Synth}
  \begin{array}{@{} l @{} c @{} c @{} l @{} l @{}}
    \lspan5{h_2 = \lbox F} \\
    \lspan5{h_{1,3,4} = \lifted\psi} \\
    f_2 = 
      \theta & i & j & \mapsto\min\langle & \psi_{ij}, \\
             & ^{^{(I_0)}} & ^{^{(J_1)}} &
                                          & \min \vtyped p {I_0} \mapsto\theta_{pj}+w_{pij}, \\
             & & &                        & \min \vtyped q {J_1} \mapsto\theta_{iq}+w'_{qji}\rangle \\
    \lspan5{f_{1,3,4} = \lifted\psi}
  \end{array}
\end{tacticbox}

\begin{equation}
  \fix\quadrants{\lifted\psi}{\lbox B}{\lifted\psi}{\lifted\psi} ~=~
    \quadrants{\psi}{B^{^{I_0J_0J_1}}_\psi\!\!\!\!}{\psi}{\psi} ~\applt~
    \psi\mapsto\quadrants{\psi}{A^{^{I_0J_1}}_\psi\!\!\!}{\psi}{\psi}
  \label{fix B}
\end{equation}

\medskip\noindent
In a similar manner, we will obtain the following:

\begin{equation}
  \fix\quadrants{\lifted\psi}{\lifted\psi}{\lbox C}{\lifted\psi} ~=~
    \quadrants{\psi}{\psi}{C^{^{I_0I_1J_0}}_\psi\!\!\!\!}{\psi} ~\applt~
    \psi\mapsto\quadrants{\psi}{\psi}{A^{^{I_1J_0}}_\psi\!\!\!}{\psi}
  \label{fix C}
\end{equation}

\begin{equation}
  \renewcommand\arraystretch{1.2}
  \begin{array}{@{}l @{} l @{\!} c @{} c @{} l @{} l@{}}
  C^{^{I_0I_1J}} =~ \big( & \psi\mapsto \\
      & \fix
      \theta & i & j & \mapsto\min\langle & \psi_{ij}, \\
           & & ^{^{(I)}} & ^{^{(J)}} &
                                          & \min \vtyped p {I_0} \mapsto\theta_{pj}+w_{pij} \rangle\big) \\
      & \lspan5{:: \big((I_0\times J)\to\R\big) \to \big((I_1\times J)\to\R\big)}
  \end{array}
\end{equation}

\medskip\noindent
And ---

\begin{equation}
  \begin{array}{@{} l @{} l @{}}
    \fix\quadrants{\lifted\psi}{\lifted\psi}{\lifted\psi}{\lbox D} ~=~ &
      \quadrants{\psi}{\psi}{\psi}{B^{^{I_1J_0J_1}}_\psi\!\!\!\!} ~\applt~
      \psi\mapsto\quadrants{\psi}{\psi}{\psi}{C^{^{I_0I_1J_1}}_\psi\!\!\!} \\
    &
       ~\applt~ \psi\mapsto\quadrants{\psi}{\psi}{\psi}{A^{^{I_1J_1}}_\psi\!\!\!}
  \end{array}
  \label{fix D}
\end{equation}

This gives the stratified version as shown in \Cref{evaluation:arbiter stratify A chain}.
The read and write regions are already encoded in the types of $A$, $B$, $C$ in 
\eqref{fix A}, \eqref{fix B}, \eqref{fix C}, and \eqref{fix D}.

\begin{figure*}
\centering
\begin{tikzpicture}[>=latex,x=7mm,y=7mm,
    every path/.style={step=1},
    every node/.style={inner sep=.5pt},
    below edge/.style={below=1mm},
    above edge/.style={above=1mm},
    block/.style={rectangle,draw,thick,fill=blue, fill opacity=0.15, inner sep=0}]
    
  \def\dx{2.1cm}
  \def\dy{1cm}
  \def\w{3mm}
    
  \draw (0,0) grid (2,2);
  \node(1) at (.5,1.5) {1};   \node(2) at (1.5,1.5) {2};
  \node(3) at (.5,.5) {3};    \node(4) at (1.5,.5) {4};
  \node(s)[block,fit={(0,1) (1,2)}] {};

  \node[inner sep=0] at (2.5,1) {\includegraphics[width=\w]{img/arrow}};
  
  \tikzset{xshift=\dx}
  
  \draw (0,0) grid (2,2);
  \node(1) at (.5,1.5) {\,1$'$};   \node(2) at (1.5,1.5) {2};
  \node(3) at (.5,.5) {3};    \node(4) at (1.5,.5) {4};
  \draw[->] (s.60) to[out=30] node[above edge] {$A$} (1);
  \node(s)[block,fit={(.05,1) (2.05,2)}] {};
  \node(t)[block,fit={(0,2.05) (1,0)}] {};

  \node[inner sep=0] at (2.5,1) {\includegraphics[width=\w]{img/arrow}};
  
  \tikzset{xshift=\dx, yshift=\dy}
 
  \draw (0,0) grid (2,2);
  \node(1) at (.5,1.5) {\,1$'$};   \node(2) at (1.5,1.5) {\,2$'$};
  \node(3) at (.5,.5) {3};    \node(4) at (1.5,.5) {4};
  \draw[->] (s.60) to[out=80,looseness=1.2] node[above edge] {$B$} (2.120);
  \node(s)[block,fit={(1,1) (2,2)}] {};

  \tikzset{yshift=-2*\dy}
 
  \draw (0,0) grid (2,2);
  \node(1) at (.5,1.5) {\,1$'$};   \node(2) at (1.5,1.5) {2};
  \node(3) at (.5,.5) {\,3$'$};    \node(4) at (1.5,.5) {4};
  \draw[->] (t.-80) to[out=-60,in=-140] node[below edge] {$C$} (3.-120);
  \node(t)[block,fit={(0,0) (1,1)}] {};

  \tikzset{yshift=\dy} % back to middle

  \node[inner sep=0] at (2.5,1) {\includegraphics[width=\w]{img/arrow}};
  
  \tikzset{xshift=\dx}
  
  \draw (0,0) grid (2,2);
  \node(1) at (.5,1.5) {\,1$'$};   \node(2) at (1.5,1.5) {\,\,2$''$};
  \node(3) at (.5,.5) {\,\,3$''$};    \node(4) at (1.5,.5) {4};
  \draw[->] (s) to[out=40,in=120] node[above edge,xshift=1mm] {$A$}  (2);
  \draw[->] (t) to[out=-60,in=-120] node[below edge] {$A$}  (3);
  \node(s)[block,fit={(1,0) (2,2)}] {};

  \node[inner sep=0] at (2.5,1) {\includegraphics[width=\w]{img/arrow}};

  \tikzset{xshift=\dx}

  \draw (0,0) grid (2,2);
  \node(1) at (.5,1.5) {\,1$'$};   \node(2) at (1.5,1.5) {\,\,2$''$};
  \node(3) at (.5,.5) {\,\,3$''$};    \node(4) at (1.5,.5) {\,4$'$};
  \draw[->] (s.-80) to[out=-40,in=-120] node[below edge] {$C$}  (4);
  \node(s)[block,fit={(0,0) (2,1)}] {};

  \node[inner sep=0] at (2.5,1) {\includegraphics[width=\w]{img/arrow}};

  \tikzset{xshift=\dx}

  \draw (0,0) grid (2,2);
  \node(1) at (.5,1.5) {\,1$'$};   \node(2) at (1.5,1.5) {\,\,2$''$};
  \node(3) at (.5,.5) {\,\,3$''$};    \node(4) at (1.5,.5) {\,\,4$''$};
  \draw[->] (s.-40) to[out=-40,in=-120] node[below edge] {$B$} (4);
  \node(s)[block,fit={(1,0) (2,1)}] {};

  \node[inner sep=0] at (2.5,1) {\includegraphics[width=\w]{img/arrow}};

  \tikzset{xshift=\dx}

  \draw (0,0) grid (2,2);
  \node(1) at (.5,1.5) {\,1$'$};   \node(2) at (1.5,1.5) {\,\,2$''$};
  \node(3) at (.5,.5) {\,\,3$''$};    \node(4) at (1.5,.5) {\,\,\,4$'''$};
  \draw[->] (s.-80) to[out=-40,in=-120] node[below edge] {$A$} (4);

\end{tikzpicture}
\caption{\label{evaluation:arbiter stratify A chain}
  Fully divide-and-conquered version of $A^{IJ}$ in the example development.}
\end{figure*}


\medskip
\hrule
\bigskip

\fi% end of \iftr


\end{document}
